%DIF 1d1
%DIF LATEXDIFF DIFFERENCE FILE
%DIF DEL old.tex   Fri Sep 25 11:28:43 2026
%DIF ADD new.tex   Fri Sep 25 11:28:43 2026
%DIF < % =====================================================================
%DIF -------
% The Silent Tax: How Typos Reshape Reasoning in Thinking LLMs
%DIF 3c2
%DIF < %  Final report -- built locally, outside the git repository.
%DIF -------
% NLP Course 2025/2026, Tel Aviv University. %DIF > 
%DIF -------
%
%DIF 5-8c4-17
%DIF < %  Provenance:
%DIF < %    Abstract, 1, 2, 4, 5,     author text, verbatim from report_text.txt
%DIF < %    6, 7, 8, 9                (re-synced 2026-09-19, second pass)
%DIF < %    3 Method                  taken from Ron's paper, unchanged
%DIF -------
% Built against the official ACL Overleaf template. Drop acl.sty, acl_natbib.bst %DIF > 
% and the ACL .cls-adjacent files next to this file (see README.md). If they are %DIF > 
% absent the document still compiles with a plain fallback preamble. %DIF > 
 %DIF > 
\documentclass[11pt]{article} %DIF > 
 %DIF > 
\IfFileExists{acl.sty}{% %DIF > 
  \usepackage[final]{acl}% %DIF > 
}{% %DIF > 
  % Fallback so the document still builds without the ACL style files. %DIF > 
  \usepackage[margin=1in]{geometry}% %DIF > 
  \usepackage{natbib}% %DIF > 
  \providecommand{\And}{\quad}% %DIF > 
  \providecommand{\AND}{\quad}% %DIF > 
%DIF -------
  %
%DIF 10-11c19
%DIF < % =====================================================================
%DIF < \documentclass[10pt,twocolumn,a4paper]{article}
%DIF -------
} %DIF > 
%DIF -------

%DIF 13c21-22
%DIF < \usepackage[margin=0.78in]{geometry}
%DIF -------
\usepackage{times} %DIF > 
\usepackage{latexsym} %DIF > 
%DIF -------
\usepackage[T1]{fontenc}
\usepackage[utf8]{inputenc}
%DIF 16-17c25
%DIF < \usepackage{times}
%DIF < \usepackage{booktabs}
%DIF -------
\usepackage{microtype} %DIF > 
%DIF -------
\usepackage{graphicx}
%DIF 19a27
\usepackage{booktabs} %DIF > 
%DIF -------
\usepackage{amsmath}
\usepackage{amssymb}
%DIF 21-29c30-31
%DIF < \usepackage{float}
%DIF < \usepackage{caption}
%DIF < \usepackage{url}
%DIF < \usepackage[hidelinks]{hyperref}
%DIF < \usepackage[round]{natbib}
%DIF < % stfloats lets a full-width figure*/table* sit at the BOTTOM of a page as well
%DIF < % as the top. Without it a two-column document can only use page tops, so wide
%DIF < % floats queue up and drift away from the text that cites them.
%DIF < \usepackage{stfloats}
%DIF -------
\usepackage{multirow} %DIF > 
\usepackage{xcolor} %DIF > 
%DIF -------

%DIF 31-42c33
%DIF < % Allow floats to take more of a page before LaTeX defers them, and make it
%DIF < % much more reluctant to flush everything onto a float-only page.
%DIF < \renewcommand{\topfraction}{0.9}
%DIF < \renewcommand{\bottomfraction}{0.7}
%DIF < \renewcommand{\textfraction}{0.08}
%DIF < \renewcommand{\floatpagefraction}{0.75}
%DIF < \renewcommand{\dbltopfraction}{0.9}
%DIF < \renewcommand{\dblfloatpagefraction}{0.75}
%DIF < \setcounter{topnumber}{3}
%DIF < \setcounter{bottomnumber}{2}
%DIF < \setcounter{totalnumber}{4}
%DIF < \setcounter{dbltopnumber}{3}
%DIF -------
\graphicspath{{figures/}} %DIF > 
%DIF -------

%DIF 44-64c35
%DIF < % Run-in paragraph headings: article's default before-skip is
%DIF < % 3.25ex plus 1ex, and flushbottom stretches that 1ex to fill the column,
%DIF < % which opens uneven gaps between consecutive \paragraph blocks.
%DIF < \makeatletter
%DIF < \renewcommand\paragraph{\@startsection{paragraph}{4}{\z@}%
%DIF <   {1.9ex \@plus 0.25ex \@minus 0.2ex}%
%DIF <   {-0.9em}%
%DIF <   {\normalfont\normalsize\bfseries}}
%DIF < \makeatother
%DIF < 
%DIF < \captionsetup{font=small,labelfont=bf,skip=3pt}
%DIF < \setlength{\bibsep}{1pt plus 0.3ex}
%DIF < \setlength{\textfloatsep}{6pt plus 1pt minus 2pt}
%DIF < \setlength{\dbltextfloatsep}{6pt plus 1pt minus 2pt}
%DIF < \setlength{\floatsep}{4pt plus 1pt minus 2pt}
%DIF < \setlength{\dblfloatsep}{4pt plus 1pt minus 2pt}
%DIF < \setlength{\intextsep}{4pt plus 1pt minus 2pt}
%DIF < \setlength{\abovecaptionskip}{3pt}
%DIF < \setlength{\belowcaptionskip}{0pt}
%DIF < \renewcommand{\arraystretch}{1.02}
%DIF < 
%DIF -------
\newcommand{\typo}[1]{\textcolor{red!75!black}{\textbf{#1}}} %DIF > 
%DIF -------
\newcommand{\pp}{\,pp}
\newcommand{\arrow}{$\rightarrow$}

\title{The Silent Tax: How Typos Reshape Reasoning in Thinking LLMs}

\author{%
  Shiran Hamami \\ 207487026 \\ \small\texttt{shiranhamami@mail.tau.ac.il}
  \And
  Dolev Abudi \\ 323096099 \\ \small\texttt{dolevabudi@mail.tau.ac.il}
  \AND
  Ron Shtricker \\ 209524552 \\ \small\texttt{ronshtricker@mail.tau.ac.il}
  \And
  Ido Azoulay \\ 212710958 \\ \small\texttt{idoazoulay1@mail.tau.ac.il}
}
%DIF < \date{}
%DIF DELETED TITLE COMMANDS FOR MARKUP
\date{}%DIFAUXCMD
%DIF PREAMBLE EXTENSION ADDED BY LATEXDIFF
%DIF UNDERLINE PREAMBLE %DIF PREAMBLE
\RequirePackage[normalem]{ulem} %DIF PREAMBLE
\RequirePackage{color}\definecolor{RED}{rgb}{1,0,0}\definecolor{BLUE}{rgb}{0,0,1} %DIF PREAMBLE
\providecommand{\DIFadd}[1]{{\protect\color{blue}\uwave{#1}}} %DIF PREAMBLE
\providecommand{\DIFdel}[1]{{\protect\color{red}\sout{#1}}} %DIF PREAMBLE
%DIF SAFE PREAMBLE %DIF PREAMBLE
\providecommand{\DIFaddbegin}{} %DIF PREAMBLE
\providecommand{\DIFaddend}{} %DIF PREAMBLE
\providecommand{\DIFdelbegin}{} %DIF PREAMBLE
\providecommand{\DIFdelend}{} %DIF PREAMBLE
\providecommand{\DIFmodbegin}{} %DIF PREAMBLE
\providecommand{\DIFmodend}{} %DIF PREAMBLE
%DIF FLOATSAFE PREAMBLE %DIF PREAMBLE
\providecommand{\DIFaddFL}[1]{\DIFadd{#1}} %DIF PREAMBLE
\providecommand{\DIFdelFL}[1]{\DIFdel{#1}} %DIF PREAMBLE
\providecommand{\DIFaddbeginFL}{} %DIF PREAMBLE
\providecommand{\DIFaddendFL}{} %DIF PREAMBLE
\providecommand{\DIFdelbeginFL}{} %DIF PREAMBLE
\providecommand{\DIFdelendFL}{} %DIF PREAMBLE
\providecommand{\DIFscaledelfig}{0.5}
%DIF HIGHLIGHTGRAPHICS PREAMBLE %DIF PREAMBLE
\RequirePackage{settobox} %DIF PREAMBLE
\RequirePackage{letltxmacro} %DIF PREAMBLE
\newsavebox{\DIFdelgraphicsbox} %DIF PREAMBLE
\newlength{\DIFdelgraphicswidth} %DIF PREAMBLE
\newlength{\DIFdelgraphicsheight} %DIF PREAMBLE
% store original definition of \includegraphics %DIF PREAMBLE
\LetLtxMacro{\DIFOincludegraphics}{\includegraphics} %DIF PREAMBLE
\providecommand{\DIFaddincludegraphics}[2][]{{\color{blue}\fbox{\DIFOincludegraphics[#1]{#2}}}} %DIF PREAMBLE
\providecommand{\DIFdelincludegraphics}[2][]{% %DIF PREAMBLE
\sbox{\DIFdelgraphicsbox}{\DIFOincludegraphics[#1]{#2}}% %DIF PREAMBLE
\settoboxwidth{\DIFdelgraphicswidth}{\DIFdelgraphicsbox} %DIF PREAMBLE
\settoboxtotalheight{\DIFdelgraphicsheight}{\DIFdelgraphicsbox} %DIF PREAMBLE
\scalebox{\DIFscaledelfig}{% %DIF PREAMBLE
\parbox[b]{\DIFdelgraphicswidth}{\usebox{\DIFdelgraphicsbox}\\[-\baselineskip] \rule{\DIFdelgraphicswidth}{0em}}\llap{\resizebox{\DIFdelgraphicswidth}{\DIFdelgraphicsheight}{% %DIF PREAMBLE
\setlength{\unitlength}{\DIFdelgraphicswidth}% %DIF PREAMBLE
\begin{picture}(1,1)% %DIF PREAMBLE
\thicklines\linethickness{2pt} %DIF PREAMBLE
{\color[rgb]{1,0,0}\put(0,0){\framebox(1,1){}}}% %DIF PREAMBLE
{\color[rgb]{1,0,0}\put(0,0){\line( 1,1){1}}}% %DIF PREAMBLE
{\color[rgb]{1,0,0}\put(0,1){\line(1,-1){1}}}% %DIF PREAMBLE
\end{picture}% %DIF PREAMBLE
}\hspace*{3pt}}} %DIF PREAMBLE
} %DIF PREAMBLE
\LetLtxMacro{\DIFOaddbegin}{\DIFaddbegin} %DIF PREAMBLE
\LetLtxMacro{\DIFOaddend}{\DIFaddend} %DIF PREAMBLE
\LetLtxMacro{\DIFOdelbegin}{\DIFdelbegin} %DIF PREAMBLE
\LetLtxMacro{\DIFOdelend}{\DIFdelend} %DIF PREAMBLE
\DeclareRobustCommand{\DIFaddbegin}{\DIFOaddbegin \let\includegraphics\DIFaddincludegraphics} %DIF PREAMBLE
\DeclareRobustCommand{\DIFaddend}{\DIFOaddend \let\includegraphics\DIFOincludegraphics} %DIF PREAMBLE
\DeclareRobustCommand{\DIFdelbegin}{\DIFOdelbegin \let\includegraphics\DIFdelincludegraphics} %DIF PREAMBLE
\DeclareRobustCommand{\DIFdelend}{\DIFOaddend \let\includegraphics\DIFOincludegraphics} %DIF PREAMBLE
\LetLtxMacro{\DIFOaddbeginFL}{\DIFaddbeginFL} %DIF PREAMBLE
\LetLtxMacro{\DIFOaddendFL}{\DIFaddendFL} %DIF PREAMBLE
\LetLtxMacro{\DIFOdelbeginFL}{\DIFdelbeginFL} %DIF PREAMBLE
\LetLtxMacro{\DIFOdelendFL}{\DIFdelendFL} %DIF PREAMBLE
\DeclareRobustCommand{\DIFaddbeginFL}{\DIFOaddbeginFL \let\includegraphics\DIFaddincludegraphics} %DIF PREAMBLE
\DeclareRobustCommand{\DIFaddendFL}{\DIFOaddendFL \let\includegraphics\DIFOincludegraphics} %DIF PREAMBLE
\DeclareRobustCommand{\DIFdelbeginFL}{\DIFOdelbeginFL \let\includegraphics\DIFdelincludegraphics} %DIF PREAMBLE
\DeclareRobustCommand{\DIFdelendFL}{\DIFOaddendFL \let\includegraphics\DIFOincludegraphics} %DIF PREAMBLE
%DIF AMSMATHULEM PREAMBLE %DIF PREAMBLE
\makeatletter %DIF PREAMBLE
\let\sout@orig\sout %DIF PREAMBLE
\renewcommand{\sout}[1]{\ifmmode\text{\sout@orig{\ensuremath{#1}}}\else\sout@orig{#1}\fi} %DIF PREAMBLE
\makeatother %DIF PREAMBLE
%DIF COLORLISTINGS PREAMBLE %DIF PREAMBLE
\RequirePackage{listings} %DIF PREAMBLE
\RequirePackage{color} %DIF PREAMBLE
\lstdefinelanguage{DIFcode}{ %DIF PREAMBLE
%DIF DIFCODE_UNDERLINE %DIF PREAMBLE
  moredelim=[il][\color{red}\sout]{\%DIF\ <\ }, %DIF PREAMBLE
  moredelim=[il][\color{blue}\uwave]{\%DIF\ >\ } %DIF PREAMBLE
} %DIF PREAMBLE
\lstdefinestyle{DIFverbatimstyle}{ %DIF PREAMBLE
	language=DIFcode, %DIF PREAMBLE
	basicstyle=\ttfamily, %DIF PREAMBLE
	columns=fullflexible, %DIF PREAMBLE
	keepspaces=true %DIF PREAMBLE
} %DIF PREAMBLE
\lstnewenvironment{DIFverbatim}[1][]{\lstset{style=DIFverbatimstyle}}{} %DIF PREAMBLE
\lstnewenvironment{DIFverbatim*}[1][]{\lstset{style=DIFverbatimstyle,showspaces=true}}{} %DIF PREAMBLE
\lstset{extendedchars=\true,inputencoding=utf8}

%DIF END PREAMBLE EXTENSION ADDED BY LATEXDIFF

\usepackage{adjustbox}
\usepackage{float}
\definecolor{DIFgreen}{rgb}{0,0.5,0.1}
\providecommand{\DIFaddB}[1]{{\protect\color{blue}\uwave{#1}}}
\providecommand{\DIFaddBFL}[1]{\DIFaddB{#1}}
\providecommand{\DIFaddG}[1]{{\protect\color{DIFgreen}\uwave{#1}}}
\providecommand{\DIFaddGFL}[1]{\DIFaddG{#1}}
\begin{document}
\maketitle
\begin{center}\fbox{\parbox{0.95\linewidth}{\small\textbf{How changes are marked.} \DIFdel{Red strikethrough}: deleted. \DIFaddB{Blue underline}: added (ACL format and writing). \DIFaddG{Green underline}: a number replaced by a different one. Figures are images and show the edited version. This file shows every change from the team's version to the edited version in the lecturer's ACL format.}}\end{center}

% ---------------------------------------------------------------------
%DIF <  ABSTRACT -- AUTHOR TEXT, verbatim.
%DIF <  ---------------------------------------------------------------------
\begin{abstract}
Reasoning models perform strongly on clean mathematical and scientific
benchmarks, but real users often write with typos. We study how these everyday
errors affect not only final-answer accuracy, but also the reasoning process
itself. \DIFdelbegin %DIFDELCMD < 

%DIFDELCMD < %%%
\DIFdelend We introduce a controlled typo generator that independently varies the
number of corrupted words and the proportion of typos that produce another
valid English word\DIFdelbegin \DIFdel{. We }\DIFdelend \DIFaddbegin \DIFaddB{, }\DIFaddend apply it to 500 questions from each of GSM8K, MATH-500 \DIFdelbegin \DIFdel{, }\DIFdelend and
ARC-Challenge, and evaluate DeepSeek-R1-Distill-Qwen-7B across all corruption
settings. \DIFdelbegin %DIFDELCMD < 

%DIFDELCMD < %%%
\DIFdelend Typos consistently reduce accuracy, \DIFdelbegin \DIFdel{with the largest effect }\DIFdelend \DIFaddbegin \DIFaddB{most }\DIFaddend on GSM8K, where \DIFdelbegin \DIFdel{performance }\DIFdelend \DIFaddbegin \DIFaddB{it }\DIFaddend falls from
91.4\% on clean questions to 60.4\% in the most corrupted setting. \DIFdelbegin \DIFdel{More importantly, typos }\DIFdelend \DIFaddbegin \DIFaddB{Typos }\DIFaddend that
produce another valid \DIFdelbegin \DIFdel{English word are consistently }\DIFdelend \DIFaddbegin \DIFaddB{word are }\DIFaddend more harmful than non-word typos\DIFdelbegin \DIFdel{. Depending }\DIFdelend \DIFaddbegin \DIFaddB{: depending }\DIFaddend on
the dataset, a real-word typo is estimated to be \DIFdelbegin \DIFdel{1.9 to 5.1 }\DIFdelend \DIFaddbegin \DIFaddG{1.3 to 5.3 }\DIFaddend times more damaging
per typo. \DIFdelbegin %DIFDELCMD < 

%DIFDELCMD < %%%
\DIFdelend The reasoning traces show that the model often notices that something
is wrong\DIFdelbegin \DIFdel{:
corrupted inputs lead to }\DIFdelend \DIFaddbegin \DIFaddB{, with }\DIFaddend longer reasoning and more uncertainty language\DIFdelbegin \DIFdel{.
However, }\DIFdelend \DIFaddbegin \DIFaddB{, but }\DIFaddend this does not
translate into more effective self-correction\DIFdelbegin \DIFdel{. The }\DIFdelend \DIFaddbegin \DIFaddB{: the }\DIFaddend model hesitates more, but
does not verify more. \DIFdelbegin %DIFDELCMD < 

%DIFDELCMD < %%%
\DIFdel{Finally, none }\DIFdelend \DIFaddbegin \DIFaddB{None }\DIFaddend of the lightweight \DIFdelbegin \DIFdel{mitigation strategies we tested consistently
improved performance. Warning the model does not improve performance, asking it
to rewrite }\DIFdelend \DIFaddbegin \DIFaddB{mitigations we tested (a typo
warning, rewriting }\DIFaddend the question first\DIFdelbegin \DIFdel{can make performance worse, and spell checking
becomes less effective as the proportion of real-word typos increases. Overall,
our results suggest that the }\DIFdelend \DIFaddbegin \DIFaddB{, or an external spell checker)
consistently improved accuracy. The }\DIFaddend main challenge is not \DIFdelbegin \DIFdel{only }\DIFdelend detecting that the
input is corrupted, but \DIFdelbegin \DIFdel{correctly recovering the }\DIFdelend \DIFaddbegin \DIFaddB{recovering its }\DIFaddend intended meaning.
\end{abstract}

% =====================================================================
\section{Introduction}
%DIF <  AUTHOR TEXT - verbatim.
% =====================================================================

Chain-of-thought prompting and reasoning-focused models have substantially
improved performance on mathematical and scientific benchmarks
\citep{wei2022,deepseekr1}. However, benchmark questions are \DIFdelbegin \DIFdel{usually }\DIFdelend carefully written,
while real \DIFdelbegin \DIFdel{user inputs often contain typos }\DIFdelend \DIFaddbegin \DIFaddB{users often make simple typing mistakes }\DIFaddend such as swapped letters,
missing characters, or nearby-key \DIFdelbegin \DIFdel{mistakes. }%DIFDELCMD < 

%DIFDELCMD < %%%
\DIFdel{This gap motivates our study: real users often make simple typing mistakes when
interacting with LLMs, so strong }\DIFdelend \DIFaddbegin \DIFaddB{errors. Strong }\DIFaddend performance on clean benchmark
text may \DIFaddbegin \DIFaddB{therefore }\DIFaddend overestimate reliability in real-world use.

Previous work has shown that typos can hurt LLM performance, often linking this
effect to subword tokenization \citep{gan2024,chai2024,zhao2026}. However, these
studies mainly focus on final-answer accuracy and do not examine what happens
inside the reasoning process.

Reasoning models make this question especially interesting. Their long reasoning
traces give them many opportunities to notice a corrupted word, infer what was
intended, and recover before producing the final answer. Do they actually do so?
\DIFdelbegin \DIFdel{We study this question using a controlled two-axis experimental design.
}\DIFdelend 

To study this, our typo generator separates two factors: how many words are
corrupted and what kind of corruption they contain. In particular, we vary the
share of typos that accidentally produce another valid English word. This allows
us to compare real-word and non-word typos while keeping the total number of
typos fixed and using the same questions in every condition\DIFdelbegin \DIFdel{- }\DIFdelend \DIFaddbegin \DIFaddB{, }\DIFaddend a comparison that\DIFdelbegin \DIFdel{previous work has }\DIFdelend \DIFaddbegin \DIFaddB{,
to our knowledge, previous LLM typo studies have }\DIFaddend not made. \DIFdelbegin %DIFDELCMD < 

%DIFDELCMD < %%%
\DIFdelend We evaluate
DeepSeek-R1-Distill-Qwen-7B \citep{deepseekr1} across the full set of conditions
on three datasets with different reasoning characteristics\DIFdelbegin \DIFdel{. Beyond
measuring final-answer accuracy, we also }\DIFdelend \DIFaddbegin \DIFaddB{, and }\DIFaddend examine the
reasoning traces themselves using word-level alignment and an LLM judge.

Our main contributions are: (i) a controlled experimental setup that
independently varies typo rate and the proportion of real-word typos across
three datasets\DIFdelbegin \DIFdel{, allowing us to directly compare real-word and non-word
corruptions while keeping the overall typo rate fixed. }\DIFdelend \DIFaddbegin \DIFaddB{; }\DIFaddend (ii) an analysis of how typos affect the reasoning process
itself, including \DIFdelbegin \DIFdel{changes in }\DIFdelend model uncertainty, whether the model notices corrupted words,
and whether it \DIFdelbegin \DIFdel{attempts
to repair them, }\DIFdelend \DIFaddbegin \DIFaddB{repairs them; }\DIFaddend and (iii) an evaluation of three simple mitigation
strategies: warning the model that the input may contain typos, asking it to
correct possible typos before solving the problem, and applying an external
\DIFdelbegin \DIFdel{spellchecker
before sending the prompt to }\DIFdelend \DIFaddbegin \DIFaddB{spell checker before }\DIFaddend the \DIFdelbegin \DIFdel{model.
}\DIFdelend \DIFaddbegin \DIFaddB{model sees the question.
}\DIFaddend 

\DIFaddbegin \DIFaddB{Our main findings are as follows. Typos cost up to 31 percentage points of
accuracy on GSM8K, and the loss comes from wrong answers rather than truncated
generations (\S\ref{sec:acc}). At a fixed number of typos, real-word typos are
more harmful than non-word ones, by 1.3--5.3$\times$ per typo
(\S\ref{sec:realword}). Contrary to our initial expectation that real-word
typos would pass unnoticed, the model flags corrupted words, and hedges, more
often as corruption grows, yet it also misreads more words and does not verify
more (\S\ref{sec:repair}, \S\ref{sec:doubt}). Finally, none of the three
mitigations reliably recovers the loss (\S\ref{sec:fixes}).
}


\DIFaddend \begin{figure}[t]
\centering
\small
\setlength{\fboxsep}{5pt}
{\scriptsize\color{red}\textbf{Old version of this box (deleted):}}\\[2pt]
\fbox{\begin{minipage}{0.93\columnwidth}
\noindent{\DIFdelFL{\textbf{Clean} (model answers 18, correct)}}

\smallskip
\noindent{\itshape\DIFdelFL{Janet's ducks lay 16 eggs per day. She eats three for breakfast
every morning and bakes muffins for her friends every day with four. She sells
the remainder at the farmers' market daily for \$2 per fresh duck egg. How much
in dollars does she make every day at the farmers' market?}}

\medskip\hrule\medskip

\noindent{\DIFdelFL{\textbf{25\% typos, 40\% of them real words} (model answers 26, wrong)}}

\smallskip
\noindent{\itshape\DIFdelFL{Janet's \textbf{duck} lay 16 eggs per \textbf{adqy}.
\textbf{Sje} \textbf{east} \textbf{rthere} for breakfast every morning and bakes
muffins for \textbf{uher} \textbf{friend} every \textbf{azy} with four. She
sells the remainder at the farmers' market daily for \$2 per \textbf{resh} duck
egg. How much in dollars does she make every day at the \textbf{faremrs'}
market?}}

\medskip
\noindent{\ttfamily\DIFdelFL{eats\arrow east \quad ducks\arrow duck \quad three\arrow rthere
  friends\arrow friend \quad fresh\arrow resh \quad her\arrow uher}}
\end{minipage}}

\medskip{\scriptsize\color{blue}\textbf{New version of this box:}}\\[2pt]
\fbox{\begin{minipage}{0.93\columnwidth}
\textbf{Clean (model answers 18, correct)}\\[2pt]
{\footnotesize Janet's ducks lay 16 eggs per day. She eats three for breakfast
every morning and bakes muffins for her friends every day with four. She sells
the remainder at the farmers' market daily for \$2 per fresh duck egg. How much
in dollars does she make every day at the farmers' market?}\\[6pt]
\textbf{$r{=}25\%$, $\rho{=}40\%$ (model answers 26, wrong)}\\[2pt]
{\footnotesize Janet's \typo{duck} lay 16 eggs per \typo{adqy}.
\typo{Sje} \typo{east} \typo{rthere} for breakfast every morning and bakes
muffins for \typo{uher} \typo{friend} every \typo{azy} with four. She
sells the remainder at the farmers' market daily for \$2 per \typo{resh} duck
egg. How much in dollars does she make every day at the \typo{faremrs'}
market?}
\end{minipage}}
\caption{A single GSM8K item under our generator. Numbers and \LaTeX{} are
protected\DIFdelbegin \DIFdel{, and }\DIFdelend \DIFaddbegin \DIFaddB{; }\DIFaddend only alphabetic words are corrupted. \DIFdelbegin \DIFdel{The two }\emph{\DIFdel{real-word}}
%DIFAUXCMD
\DIFdelend \DIFaddbegin \DIFaddB{Three of the ten }\DIFaddend typos \DIFaddbegin \DIFaddB{are
}\emph{\DIFaddB{real words}} \DIFaddend (\texttt{eats\arrow east}, \texttt{ducks\arrow duck}\DIFaddbegin \DIFaddB{,
}\texttt{\DIFaddB{friends\arrow friend}}\DIFaddend ) \DIFaddbegin \DIFaddB{and }\DIFaddend are invisible to a spell checker. The model
\DIFdelbegin \DIFdel{loses the ``three eaten'' quantity and }\DIFdelend answers 26 instead of 18.}
\label{fig:example}
\end{figure}


% =====================================================================
\section{Background and Related Work}
%DIF <  AUTHOR TEXT - verbatim.
% =====================================================================

\paragraph{Sensitivity to character-level typos.}
\citet{belinkov2018} show that both synthetic and natural character-level noise
substantially degrade neural machine translation, and that robustness to one
type of noise does not necessarily transfer to others. Later work studied
similar vulnerabilities in adversarial settings. HotFlip \citep{ebrahimi2018}
shows that a few white-box character changes can alter a classifier's
prediction, while TextFooler \citep{jin2020} generates black-box adversarial
examples through word substitutions. PromptRobust \citep{zhu2024} shows that
small \DIFaddbegin \DIFaddB{character- and word-level }\DIFaddend changes to task prompts can substantially
reduce LLM performance\DIFdelbegin \DIFdel{,
including changes at both the character and word levels}\DIFdelend . \citet{pruthi2019} defend against adversarial
misspellings by placing a robust word-recognition model before the classifier\DIFdelbegin \DIFdel{. In one experiment,
an adversarial character attack
reduced BERT accuracy from 90.3\% to 45.8\%, while their defense recovered it to
75\%. This approach is }\DIFdelend \DIFaddbegin \DIFaddB{,
an approach }\DIFaddend closely related to the spell-checking mitigation we evaluate.

\paragraph{Typos and modern LLMs.}
\DIFdelbegin \DIFdel{Recent studies show that typos can substantially reduce LLM performance.
}\DIFdelend \citet{gan2024} use an optimized adversarial search to inject typos into
reasoning benchmarks such as GSM8K, BBH, and MMLU, and show that model accuracy
drops substantially as the number of character-level errors increases.
\citet{chai2024} show that character-level typos can hurt model performance
partly because they change how words are split into subwords.
\citet{zhao2026} study typo robustness across several languages and 18
open-source models\DIFdelbegin \DIFdel{. They }\DIFdelend \DIFaddbegin \DIFaddB{, and }\DIFaddend find that typos hurt \DIFdelbegin \DIFdel{performance }\DIFdelend most on tasks that require
generation or more complex reasoning.

These studies measure final-answer accuracy rather than the reasoning process
itself, and they vary typos by generation mechanism rather than by whether the
corrupted form is itself a valid word. That distinction matters. A typo that
creates a non-word, such as \texttt{triangle\arrow trianlge}, mainly affects
tokenization, while a typo that creates another valid word, such as
\texttt{sum\arrow sun}, leaves tokenization intact and changes the meaning of
the sentence. Our study separates these cases and examines how the model
responds to each during reasoning.

\paragraph{Reasoning traces and self-correction.}
Reasoning models such as DeepSeek-R1 \citep{deepseekr1} produce long reasoning
traces before giving a final answer\DIFdelbegin \DIFdel{, allowing us to examine how they respond to
errors during reasoning. However, these }\DIFdelend \DIFaddbegin \DIFaddB{. These }\DIFaddend traces do not always fully reflect
the model's internal reasoning \citep{turpin2023}, and \DIFdelbegin \DIFdel{prior work shows that }\DIFdelend LLMs often struggle to
correct their own reasoning without external feedback, sometimes even degrading
their answers in the process \citep{huang2024}. \DIFdelbegin %DIFDELCMD < 

%DIFDELCMD < %%%
\DIFdelend Our setting asks a related but
different question: can a long reasoning process help the model deal with errors
that are already present in the input? A typo may be detectable from context,
giving the model an opportunity to notice it, infer the intended word, and
continue reasoning correctly.
\DIFdelbegin \DIFdel{This allows us to
study whether additional reasoning actually helps models recover from noisy
input.
}\DIFdelend 

% =====================================================================
\section{Method}
%DIF <  AUTHOR TEXT - verbatim.
% =====================================================================

\subsection{Controlled typo generation}

We generate corrupted versions of each question by modifying only its
natural-language text. Numbers, \LaTeX{} expressions, and \LaTeX{} command
names are protected and never modified. Candidate words must consist only of
ASCII letters, and we exclude single-character words and a small set of common
function words.

The generator varies two factors independently. The typo rate $r$ determines
how many eligible words are corrupted, while the real-word ratio $\rho$
determines the target fraction of those corruptions that should result in
another valid English word. For example, \texttt{eats}~\arrow~\texttt{east}
is a real-word typo, while \texttt{ducks}~\arrow~\texttt{dcks} is a non-word
typo. Words are selected for corruption before $\rho$ is applied, so for a
fixed typo rate the number of corrupted words is identical across all values
of $\rho$.

Each corrupted word receives a character-level edit sampled with weights 0.40
adjacent-key substitution, 0.20 deletion, 0.20 insertion, and 0.20
transposition. With probability 0.10, a second edit is applied to the same
word.

For each question, we convert $\rho$ into a target number of real-word typos,
using probabilistic rounding when the target is not an integer. We then process
the selected words in random order. While the target has not been reached, the
generator searches for valid real-word variants using the NLTK English lexicon.
If none exists for a selected word, that word receives a non-word typo and the
generator continues trying to satisfy the target with the remaining words. Any
real word produced accidentally is counted toward the target. Because the
selected words are not chosen based on whether a real-word variant exists, the
requested ratio cannot always be reached exactly\DIFdelbegin \DIFdel{. In these cases, the generator }\DIFdelend \DIFaddbegin \DIFaddB{; the generator then }\DIFaddend gets as
close to the target as the selected words allow.
Table~\ref{tab:datastats} reports the achieved ratios.

For each benchmark, we create a clean condition that is identical to the
original dataset, together with nine corrupted conditions. These combine typo
rates $r \in \{25, 50, 75\}\%$ with real-word ratios
$\rho \in \{10, 40, 70\}\%$. The generator is seeded consistently across
conditions, and the same questions are used in every condition, allowing direct
paired comparisons. The resulting dataset variants are available on the Hugging
Face Hub.

\subsection{Datasets}

We use 500 questions from each of three benchmarks: GSM8K \citep{cobbe2021},
MATH-500 \citep{hendrycks2021,lightman2024}, and ARC-Challenge
\citep{clark2018}. \DIFdelbegin \DIFdel{We chose these datasets because they }\DIFdelend \DIFaddbegin \DIFaddB{They }\DIFaddend differ in how much of the information needed to solve a
question is expressed in natural-language text. \DIFdelbegin %DIFDELCMD < 

%DIFDELCMD < %%%
\DIFdelend GSM8K contains arithmetic word
problems \DIFdelbegin \DIFdel{in which most of the problem }\DIFdelend \DIFaddbegin \DIFaddB{whose }\DIFaddend structure is described \DIFdelbegin \DIFdel{in prose. Its questions are also relatively long, so more words
are available for corruption}\DIFdelend \DIFaddbegin \DIFaddB{almost entirely in prose, and its
questions are relatively long}\DIFaddend . MATH-500 contains \DIFdelbegin \DIFdel{competition mathematics
problems
that are shorter and contain }\DIFdelend \DIFaddbegin \DIFaddB{shorter competition problems
with }\DIFaddend substantial \LaTeX{} content, which \DIFdelbegin \DIFdel{is
protected by our generator }\DIFdelend \DIFaddbegin \DIFaddB{our generator protects}\DIFaddend . ARC-Challenge
contains short\DIFdelbegin \DIFdel{er}\DIFdelend  multiple-choice science questions, \DIFdelbegin \DIFdel{where the answer options provide additional information that
}\DIFdelend \DIFaddbegin \DIFaddB{whose answer options }\DIFaddend may help
the model interpret a corrupted question.

\DIFdelbegin \DIFdel{These differences mean that }\DIFdelend \DIFaddbegin \DIFaddB{As a result, }\DIFaddend the same typo rate does not produce the same absolute number of
typos across datasets\DIFdelbegin \DIFdel{. }\DIFdelend \DIFaddbegin \DIFaddB{: }\DIFaddend GSM8K receives the most corruption per question
\DIFdelbegin \DIFdel{, while MATH-500 and ARC-Challenge receive fewer corruptions
because they contain fewer eligible words}\DIFdelend \DIFaddbegin \DIFaddB{(Table~\ref{tab:datastats})}\DIFaddend . We therefore compare corruption levels within each
dataset rather than \DIFdelbegin \DIFdel{treating a given typo rate as the same
absolute amount of noise across all three benchmarks.
Table~\ref{tab:datastats} reports the resulting corruption statistics.
}\DIFdelend \DIFaddbegin \DIFaddB{across benchmarks.
}\DIFaddend 


\begin{table}[t]
\centering
\footnotesize
\setlength{\tabcolsep}{3.5pt}
\resizebox{\columnwidth}{!}{\begin{tabular}{lrrr}
\toprule
 & GSM8K & MATH-500 & ARC-C \\
\midrule
questions & 500 & 500 & 500 \\
mean words per question & 45.8 & 31.5 & 23.5 \\
median words per question & 43 & 26 & 20 \\
eligible words per question & 30.6 & 15.7 & 16.6 \\
\midrule
typos per question, $r{=}25\%$ & 7.6 & 4.0 & 4.2 \\
typos per question, $r{=}50\%$ & 15.3 & 7.9 & 8.3 \\
typos per question, $r{=}75\%$ & 22.9 & 11.8 & 12.4 \\
\midrule
achieved real-word ratio, $\rho{=}10\%$ & 0.100 & 0.099 & 0.102 \\
achieved real-word ratio, $\rho{=}40\%$ & 0.399 & 0.397 & 0.401 \\
achieved real-word ratio, $\rho{=}70\%$ & 0.698 & 0.681 & 0.682 \\
\midrule
total corrupted words & 68{,}739 & 35{,}685 & 37{,}248 \\
\bottomrule
\end{tabular}}
\caption{Dataset and corruption statistics over the nine corrupted conditions
($r\in\{25,50,75\}\%$, $\rho\in\{10,40,70\}\%$), first 500 questions \DIFdelbegin \DIFdel{per
benchmark}\DIFdelend \DIFaddbegin \DIFaddG{of GSM8K
and ARC-Challenge and all of MATH-500}\DIFaddend . \emph{Eligible words} are those the
generator may corrupt, after protected spans, single-character words and
function words are removed; the typo rate applies to them, and real-word ratios
are pooled over all corrupted words.}
\label{tab:datastats}
\end{table}


% =====================================================================
\section{Experimental Setup}
%DIF <  AUTHOR TEXT - verbatim.
\DIFaddbegin \label{sec:setup}
\DIFaddend % =====================================================================

\paragraph{Models.}
We use DeepSeek-R1-Distill-Qwen-7B \citep{deepseekr1}\DIFdelbegin \DIFdel{as our main model and
access it }\DIFdelend \DIFaddbegin \DIFaddB{, accessed }\DIFaddend through the
Hugging Face API\DIFdelbegin \DIFdel{. We run it }\DIFdelend \DIFaddbegin \DIFaddB{, }\DIFaddend with temperature 0.6 and top-$p$ 0.95 \DIFdelbegin \DIFdel{, using the default settings from the model card, and generate one
response per question. We use the recommended decoding settings to better
}\DIFdelend \DIFaddbegin \DIFaddB{(the model-card
defaults, to }\DIFaddend reflect how a typical user would \DIFdelbegin \DIFdel{interact with the model, rather than tuning
them specifically for our evaluation. For each dataset, we set the }\DIFdelend \DIFaddbegin \DIFaddB{run it) and one response per
question. The }\DIFaddend maximum generation length \DIFdelbegin \DIFdel{to a value that }\DIFdelend \DIFaddbegin \DIFaddG{is 20}{\DIFaddB{,}}\DIFaddB{000 tokens for GSM8K and
17}{\DIFaddB{,}}\DIFaddB{000 for MATH-500 and ARC-Challenge, which }\DIFaddend we found sufficient for the model
to \DIFdelbegin \DIFdel{produce
complete answers for a large enough portion of the questions. }\footnote{\DIFdel{In
practice this is 20}%DIFDELCMD < {%%%
\DIFdel{,}%DIFDELCMD < }%%%
\DIFdel{000 generated tokens for GSM8K and 17}%DIFDELCMD < {%%%
\DIFdel{,}%DIFDELCMD < }%%%
\DIFdel{000 for MATH-500
and ARC-Challenge.}}
%DIFAUXCMD
\addtocounter{footnote}{-1}%DIFAUXCMD
%DIFDELCMD < 

%DIFDELCMD < %%%
\DIFdelend \DIFaddbegin \DIFaddB{finish on a large majority of questions. }\DIFaddend For the LLM-based evaluation \DIFdelbegin \DIFdel{, }\DIFdelend we use
Llama-3.3-70B-Instruct \citep{llama3}, also through the Hugging Face API, with
temperature 0.

\paragraph{Prompts.}
For math problems, we provide the question followed by the instruction:
``Please reason step by step, and put your final answer within
\texttt{\textbackslash boxed\{\}}.'' For multiple-choice questions, we present
the four answer options with shuffled labels and ask the model to return the
final answer as a boxed letter. Typos are added only to the question itself,
while the instruction remains unchanged and typo-free.

\paragraph{Scoring.}
We extract the final answer only from the text produced after the closing
\texttt{</think>} tag. Responses that reach the token limit without producing a
final answer are marked as unanswered. GSM8K is evaluated using normalized
numeric matching, MATH-500 using symbolic equivalence with
\texttt{math\_verify}, and ARC using exact answer-letter matching. We report
accuracy over answered questions \DIFdelbegin \DIFdel{and separately report }\DIFdelend \DIFaddbegin \DIFaddB{together with }\DIFaddend the number of answered
responses\DIFaddbegin \DIFaddB{, and check the effect of truncation with strict accuracy, which
counts unanswered responses as wrong}\DIFaddend .

\paragraph{Statistics.}
Since each condition contains the same 500 questions, we use paired \DIFdelbegin \DIFdel{statistical
}\DIFdelend tests. We
\DIFdelbegin \DIFdel{report the number of }\DIFdelend \DIFaddbegin \DIFaddB{count the }\DIFaddend questions that change from correct to incorrect and from incorrect to
correct, and \DIFdelbegin \DIFdel{use a continuity-corrected McNemar test to
}\DIFdelend compare these paired outcomes \DIFdelbegin \DIFdel{. This analysis helps separate }\DIFdelend \DIFaddbegin \DIFaddB{with a continuity-corrected McNemar
test, which separates }\DIFaddend the effect of typos from differences in question
difficulty\DIFdelbegin \DIFdel{, since each comparison is made on the same set of questions}\DIFdelend \DIFaddbegin \DIFaddB{. Tables report uncorrected $p$-values; in the text we also state
which results survive a Holm correction over the nine tests of each family
within a dataset}\DIFaddend .

\paragraph{Reasoning trace analysis.}
We analyze the \DIFdelbegin \DIFdel{model's }\DIFdelend reasoning traces in three ways. First, \DIFdelbegin \DIFdel{we use }\DIFdelend word-level alignment
\DIFdelbegin \DIFdel{to track }\DIFdelend \DIFaddbegin \DIFaddB{tracks }\DIFaddend how the model handles each corrupted word\DIFdelbegin \DIFdel{. We check }\DIFdelend \DIFaddbegin \DIFaddB{: }\DIFaddend whether the trace uses the
original word, the corrupted \DIFdelbegin \DIFdel{version, or an explicit
correction, and classify }\DIFdelend \DIFaddbegin \DIFaddB{form, or both, classifying }\DIFaddend each typo as silently
fixed, flagged, misread, or not used \DIFdelbegin \DIFdel{.
}%DIFDELCMD < 

%DIFDELCMD < %%%
\DIFdelend \DIFaddbegin \DIFaddB{(\S\ref{sec:repair}). }\DIFaddend Second, we \DIFdelbegin \DIFdel{measure self-doubt directly from the reasoning traces. We define sets
of words and phrases associated with uncertainty and self-correction, and count
how often they appear in each trace. Since traces differ in length, we normalize
these counts }\DIFdelend \DIFaddbegin \DIFaddB{count
curated uncertainty and second-guessing expressions, normalized }\DIFaddend by trace length
\DIFdelbegin \DIFdel{. We then compare these measures across typo
conditions and between correct and incorrect answers to examine whether higher
uncertainty is associated with model failure.
}%DIFDELCMD < 

%DIFDELCMD < %%%
\DIFdel{Third, we use }\DIFdelend \DIFaddbegin \DIFaddB{(\S\ref{sec:doubt}). Third, }\DIFaddend an LLM judge \DIFdelbegin \DIFdel{. The judge receives the clean
question, the corrupted question , the list of changed words}\DIFdelend \DIFaddbegin \DIFaddB{receives the list of corrupted words
with their originals, the question as shown to the model}\DIFaddend , and the \DIFdelbegin \DIFdel{model's
reasoning
trace }\DIFdelend \DIFaddbegin \DIFaddG{reasoning
trace (its first and last 2}{\DIFaddB{,}}\DIFaddB{000 characters)}\DIFaddend . Using a \DIFdelbegin \DIFdel{predefined evaluation prompt}\DIFdelend \DIFaddbegin \DIFaddB{fixed rubric}\DIFaddend , it assigns
a \DIFdelbegin \DIFdel{self-doubt
}\DIFdelend \DIFaddbegin \emph{\DIFaddB{self-doubt}} \DIFaddend score from 0 \DIFaddbegin \DIFaddB{(a clear, straight-line solution) }\DIFaddend to 10
\DIFdelbegin \DIFdel{, which measures the ``amount of unambiguity'' and a repair
behavior }\DIFdelend \DIFaddbegin \DIFaddB{(pervasive uncertainty, reversals or looping), and a }\emph{\DIFaddB{repair}} \DIFaddend score from 0
\DIFaddbegin \DIFaddB{(the intended meaning of the question is fully understood) }\DIFaddend to 5 \DIFdelbegin \DIFdel{which measures how much the model
understands the real meaningof the question. }\DIFdelend \DIFaddbegin \DIFaddB{(the model
reasons from the wrong meaning). Higher is worse on both scales. The judge
scores answered traces only.
}\DIFaddend 

% =====================================================================
\section{Results}
%DIF <  AUTHOR TEXT - verbatim.
% =====================================================================

\subsection{Effect of typos on accuracy}
\DIFaddbegin \label{sec:acc}
\DIFaddend 


\DIFaddbegin \DIFdelbegin %DIFDELCMD < \begin{figure*}[tbp]
%DIFDELCMD < %%%
\DIFdelendFL \DIFaddbeginFL \begin{figure*}[t]
\DIFaddendFL \centering
\includegraphics[width=\textwidth]{figures/fig_accuracy.pdf}
\caption{\textbf{The accuracy tax.} \emph{(a)} answered-only accuracy against
the typo rate; the line is the mean over the three real-word ratios and the band
spans them. \emph{(b)} the same as a change from clean, in percentage points.
\emph{(c)} mean generated tokens relative to clean. Accuracy falls with the typo
rate on all three datasets and falls fastest on GSM8K, where the text carries
the whole problem. \DIFdelbeginFL \DIFdelFL{Reasoning }\DIFdelendFL \DIFaddbeginFL \DIFaddBFL{On average, reasoning }\DIFaddendFL gets \emph{longer} under typos\DIFdelbeginFL \DIFdelFL{everywhere}\DIFdelendFL , so the
loss is not the model giving up early.}
\label{fig:accuracy}
\end{figure*}


\DIFaddend Figure~\ref{fig:accuracy} and Table~\ref{tab:grid} in the appendix show the full
results. On GSM8K, answered-only accuracy drops from 91.4\% on clean questions
to 60.4\% in the most corrupted condition ($r{=}75$, $\rho{=}70$), a decrease of
31.0 \DIFdelbegin \DIFdel{percent. ARC shows a smaller drop }\DIFdelend \DIFaddbegin \DIFaddB{percentage points. ARC drops }\DIFaddend from 84.7\% to 74.7\%, \DIFdelbegin \DIFdel{while
}\DIFdelend \DIFaddbegin \DIFaddB{and }\DIFaddend MATH-500 \DIFdelbegin \DIFdel{drops }\DIFdelend from
95.2\% to 86.9\%. \DIFdelbegin \DIFdel{Overall, accuracy decreases as the typo
rate increases across all three datasets (Figure~\ref{fig:accuracy}) }\DIFdelend \DIFaddbegin \DIFaddG{On all three datasets accuracy is lowest at the highest typo
rate, although on ARC the mildest rate ($r{=}25$) is still on par with clean}\DIFaddend .

The paired analysis shows the same trend. In the most corrupted GSM8K condition,
164 questions change from correct to incorrect, while only 14 change from
incorrect to correct, out of 480 questions answered in both conditions
($p < 10^{-6}$). ARC shows 75 correct-to-incorrect versus 26
incorrect-to-correct changes, and MATH-500 shows 46 versus 7, both
\DIFdelbegin \DIFdel{statistically
significant. }\DIFdelend \DIFaddbegin \DIFaddB{significant. After Holm correction, all nine }\DIFaddend GSM8K \DIFdelbegin \DIFdel{is affected even at the mildest typo condition, while
significanteffects appear only at higher typo rates for }\DIFdelend \DIFaddbegin \DIFaddB{conditions remain
significant, against four on }\DIFaddend MATH-500 \DIFdelbegin \DIFdel{and ARC
(Table~\ref{tab:grid}). }%DIFDELCMD < 

%DIFDELCMD < %%%
\DIFdelend \DIFaddbegin \DIFaddB{(all at $r \geq 50$) and two on ARC (both
at $r{=}75$). }\DIFaddend One possible explanation \DIFdelbegin \DIFdel{for these differences }\DIFdelend is how much each task depends on the
wording of the question\DIFdelbegin \DIFdel{. }\DIFdelend \DIFaddbegin \DIFaddB{: }\DIFaddend GSM8K problems rely \DIFdelbegin \DIFdel{heavily }\DIFdelend on the text to describe the
quantities and relationships needed for the solution\DIFdelbegin \DIFdel{. ARCprovides
multiple-choice options that }\DIFdelend \DIFaddbegin \DIFaddB{, ARC's answer options }\DIFaddend may
help the model recover some information, \DIFdelbegin \DIFdel{while
}\DIFdelend \DIFaddbegin \DIFaddB{and }\DIFaddend much of the \DIFdelbegin \DIFdel{important mathematical content
in }\DIFdelend \DIFaddbegin \DIFaddB{mathematical content
of }\DIFaddend MATH-500 is written in protected \LaTeX{}\DIFdelbegin \DIFdel{and is therefore not affected by our typo generator}\DIFdelend .

\DIFdelbegin \DIFdel{Typos also lead to longer reasoning traces across all three datasets, while
the model still produces final answers in most cases. On }\DIFdelend \DIFaddbegin \DIFaddB{Truncation does not explain the loss. Counting unanswered responses as wrong,
}\DIFaddend GSM8K \DIFdelbegin \DIFdel{, for example, the average reasoninglength increases }\DIFdelend \DIFaddbegin \DIFaddG{accuracy falls from 91.4\% to 58.0\%; of this strict drop at $r{=}75$,
$\rho{=}70$, 93\% comes from answered-but-wrong responses and only 7\% from
truncation, and on MATH-500 and ARC truncation contributes nothing. Typos also
lengthen the reasoning: on GSM8K the mean trace grows }\DIFaddend from 1,582 tokens on clean
questions to 2,\DIFaddbegin \DIFaddB{306 in the most corrupted condition (and up to 2,}\DIFaddend 594 \DIFdelbegin \DIFdel{tokens }\DIFdelend at
$r{=}75$, $\rho{=}40$\DIFaddbegin \DIFaddB{)}\DIFaddend , with smaller increases on MATH-500 and ARC. \DIFdelbegin \DIFdel{This suggests that the drop in accuracy is driven mainly by incorrect
answers rather than incomplete generations. In fact, wrong answers tend to have longer reasoning }\DIFdelend \DIFaddbegin \DIFaddB{Wrong
answers have longer }\DIFaddend traces than correct ones \DIFdelbegin \DIFdel{, indicating that }\DIFdelend \DIFaddbegin \DIFaddB{in every configuration, so
}\DIFaddend additional reasoning is \DIFdelbegin \DIFdel{more closely }\DIFdelend associated with difficulty \DIFaddbegin \DIFaddB{rather }\DIFaddend than with successful
recovery (Figure~\ref{fig:accuracy}c; Table~\ref{tab:support}).

\DIFdelbegin %DIFDELCMD < \FLOATPLACEHOLDER{fig:accuracy}
%DIFDELCMD < 

%DIFDELCMD < %%%
\DIFdelend \subsection{Real-word typos are more harmful}
\DIFaddbegin \label{sec:realword}
\DIFaddend 


\DIFdelbegin %DIFDELCMD < \begin{figure}[!tbp]
%DIFDELCMD < %%%
\DIFdelendFL \DIFaddbeginFL \begin{figure}[t]
\DIFaddendFL \centering
\includegraphics[width=\columnwidth]{figures/fig_realword.pdf}
\caption{\textbf{The real-word effect.} \emph{(a)} multiplicative loss in the
odds of a correct answer per additional typo, from a logistic regression fitting
real-word and non-word counts separately. \emph{(b)} the controlled paired
comparison: the same questions with the same number of typos, changing only the
real-word share from $\rho{=}10$ to $\rho{=}70$. $^{*}$ marks McNemar
$p<0.05$ \DIFaddbeginFL \DIFaddBFL{(uncorrected)}\DIFaddendFL .}
\label{fig:realword}
\end{figure}


Because typo rate and real-word ratio are varied independently, we can compare
real-word and non-word typos while controlling for the total number of typos.
For each dataset, we fit a logistic regression \DIFdelbegin \DIFdel{using the }\DIFdelend \DIFaddbegin \DIFaddB{that predicts whether the final
answer is correct from the }\DIFaddend number of real-word and non-word typos
\DIFdelbegin \DIFdel{to predict whether the final answer is correct
}\DIFdelend (Figure~\ref{fig:realword}a). \DIFdelbegin %DIFDELCMD < 

%DIFDELCMD < %%%
\DIFdel{The results show that an }\DIFdelend \DIFaddbegin \DIFaddB{An }\DIFaddend additional real-word typo has a stronger
negative association with correctness than an additional non-word typo\DIFdelbegin \DIFdel{. Based on the
estimated odds multipliers from the logistic regression, the }\DIFdelend \DIFaddbegin \DIFaddB{: the
}\DIFaddend relative per-typo loss is about 1.9$\times$ larger on GSM8K, 1.3$\times$ on
MATH-500, and 5.3$\times$ on ARC (Table~\ref{tab:realword-logit}).

The paired comparisons support this pattern, although the effect depends on the
overall typo rate. On GSM8K at a 75\% typo rate, increasing the real-word share
from 10\% to 70\% lowers accuracy by 14.1 percentage points, even though the
same questions contain the same total number of typos. The difference is also
significant at a 50\% typo rate \DIFaddbegin \DIFaddB{(both after Holm correction)}\DIFaddend , but not at 25\%.
\DIFdelbegin \DIFdel{This suggests that when
corruption is relatively }\DIFdelend \DIFaddbegin \DIFaddB{When corruption is }\DIFaddend mild, the model can often handle both \DIFdelbegin \DIFdel{real-word and
non-word typos, while the
difference between the two }\DIFdelend \DIFaddbegin \DIFaddB{kinds of typo, and the
difference }\DIFaddend becomes more pronounced as the overall typo rate increases
(Figure~\ref{fig:realword}b; Table~\ref{tab:realword-paired}). \DIFdelbegin %DIFDELCMD < 

%DIFDELCMD < %%%
\DIFdel{ARCshows a significant real-word effect at the highest typo rate, while
}\DIFdelend \DIFaddbegin \DIFaddG{On ARC, the
effect is significant at $r{=}75$ ($-6.7$\pp{}, surviving correction); two
contrasts at $r{=}25$ reach $p<0.05$ but do not survive correction, and none at
$r{=}50$ does. }\DIFaddend MATH-500 \DIFdelbegin \DIFdel{does not show a }\DIFdelend \DIFaddbegin \DIFaddB{shows no }\DIFaddend significant difference at any typo rate\DIFdelbegin \DIFdel{. This is
consistent with the }\DIFdelend \DIFaddbegin \DIFaddB{,
consistent with its }\DIFaddend smaller number of corrupted words\DIFdelbegin \DIFdel{in MATH-500 questions,
since much of their mathematical content is protected from modification.
All
nine contrasts, for all three datasets, are in
Table~\ref{tab:realword-paired}.
}\DIFdelend \DIFaddbegin \DIFaddB{.
}\DIFaddend 

Overall, \DIFdelbegin \DIFdel{the results suggest that }\DIFdelend real-word typos \DIFdelbegin \DIFdel{can be }\DIFdelend \DIFaddbegin \DIFaddB{appear }\DIFaddend especially difficult for the model. Unlike
non-word typos, which \DIFdelbegin \DIFdel{often }\DIFdelend produce visibly unusual word forms, \DIFdelbegin \DIFdel{real-word typos }\DIFdelend \DIFaddbegin \DIFaddB{they }\DIFaddend result in valid
words and can \DIFdelbegin \DIFdel{therefore }\DIFdelend change the meaning of the input without creating an obvious
tokenization problem.

\subsection{Repair behavio\DIFdelbegin \DIFdel{u}\DIFdelend r : noticed but not repaired}
\DIFaddbegin \label{sec:repair}
\DIFaddend 


\DIFdelbegin %DIFDELCMD < \begin{figure*}[!tbp]
%DIFDELCMD < %%%
\DIFdelendFL \DIFaddbeginFL \begin{figure*}[t]
\DIFaddendFL \centering
\includegraphics[width=\textwidth]{figures/fig_repair.pdf}
\caption{\DIFdelbeginFL \textbf{\DIFdelFL{Word-level repair behaviour.}} %DIFAUXCMD
\DIFdelendFL \DIFaddbeginFL \textbf{\DIFaddBFL{Word-level repair behavior.}} \DIFaddendFL How every corrupted word is
handled in the reasoning trace. \emph{(a)} against the typo rate, pooled over
real-word ratios. \emph{(b)} against the real-word ratio, pooled over typo
rates. \DIFdelbeginFL \DIFdelFL{Section~5.3 claims a trend on both axes and both are present: silent
}\DIFdelendFL \DIFaddbeginFL \DIFaddBFL{Silent }\DIFaddendFL fixing gives way to explicit flagging and to \DIFdelbeginFL \DIFdelFL{outright }\DIFdelendFL misreading as either
\DIFdelbeginFL \DIFdelFL{knob is turned up}\DIFdelendFL \DIFaddbeginFL \DIFaddBFL{the typo rate or the real-word ratio increases}\DIFaddendFL , while the share \DIFaddbeginFL \DIFaddBFL{of words }\DIFaddendFL never
referred to \DIFdelbeginFL \DIFdelFL{at all }\DIFdelendFL stays roughly constant. Per-configuration values are in
Table~\ref{tab:repair}.}
\label{fig:repair}
\end{figure*}


We recover each corrupted word by comparing the clean and corrupted versions of
the question, and \DIFdelbegin \DIFdel{then }\DIFdelend classify how the model handles it in the reasoning trace:
\DIFdelbegin %DIFDELCMD < 

%DIFDELCMD < \begin{itemize}\setlength{\itemsep}{0pt}\setlength{\parskip}{0pt}
%DIFDELCMD <   \item %%%
\DIFdelend \textbf{silent fix}\DIFdelbegin \DIFdel{: }\DIFdelend \DIFaddbegin \DIFaddB{, }\DIFaddend only the original word appears \DIFdelbegin \DIFdel{, suggesting that }\DIFdelend \DIFaddbegin \DIFaddB{(}\DIFaddend the model recovered the
intended word without \DIFdelbegin \DIFdel{explicitly }\DIFdelend mentioning the typo\DIFdelbegin %DIFDELCMD < \item %%%
\DIFdelend \DIFaddbegin \DIFaddB{); }\DIFaddend \textbf{flagged}\DIFdelbegin \DIFdel{: }\DIFdelend \DIFaddbegin \DIFaddB{, }\DIFaddend both the original
and \DIFdelbegin \DIFdel{corrupted forms appear ,
  indicating that }\DIFdelend \DIFaddbegin \DIFaddB{the corrupted form appear (}\DIFaddend the model noticed \DIFdelbegin \DIFdel{the typo and corrected it
  }%DIFDELCMD < \item %%%
\DIFdelend \DIFaddbegin \DIFaddB{and corrected the typo);
}\DIFaddend \textbf{misread}\DIFdelbegin \DIFdel{: }\DIFdelend \DIFaddbegin \DIFaddB{, }\DIFaddend only the corrupted form appears \DIFdelbegin \DIFdel{, suggesting that the model interpreted }\DIFdelend \DIFaddbegin \DIFaddB{(the model took }\DIFaddend it as the
wrong word\DIFdelbegin %DIFDELCMD < \item %%%
\DIFdelend \DIFaddbegin \DIFaddB{); and }\DIFaddend \textbf{not used}\DIFdelbegin \DIFdel{: }\DIFdelend \DIFaddbegin \DIFaddB{, }\DIFaddend neither form appears\DIFdelbegin \DIFdel{in the reasoning trace
}%DIFDELCMD < \end{itemize}
%DIFDELCMD < %%%
\DIFdelend \DIFaddbegin \DIFaddB{.
}\DIFaddend 

We originally expected real-word typos to go unnoticed more often, while
non-word typos would be easier to detect and correct explicitly. The results
only partly support this prediction. \DIFdelbegin %DIFDELCMD < 

%DIFDELCMD < %%%
\DIFdel{Both a higher typo rate and a higher real-word ratio change the way the model
handles corrupted words. }\DIFdelend As the typo rate increases from 25\% to
75\%, silent fixing decreases \DIFdelbegin \DIFdel{across }\DIFdelend \DIFaddbegin \DIFaddB{on }\DIFaddend all three datasets\DIFdelbegin \DIFdel{. On }\DIFdelend \DIFaddbegin \DIFaddB{; on }\DIFaddend GSM8K, for example, it
falls from 73.3\% to 55.8\% at $\rho{=}10$. \DIFdelbegin \DIFdel{Similar decreases appear on ARC and MATH-500.
}\DIFdelend At the same time, flagging and
misreading become more common\DIFdelbegin \DIFdel{, and misreading roughly doubles on all three datasets }\DIFdelend \DIFaddbegin \DIFaddG{: pooled over real-word ratios, misreading rises
from 4.0\% to 8.3\% on GSM8K and from 1.5\% to 2.9\% on MATH-500, and less
steeply on ARC, from 3.5\% to 5.1\% }\DIFaddend (Figure~\ref{fig:repair};
Table~\ref{tab:repair}).

Increasing the real-word ratio \DIFdelbegin \DIFdel{produces }\DIFdelend \DIFaddbegin \DIFaddB{has }\DIFaddend a similar effect. Pooled across typo rates on
GSM8K, silent fixing falls from 62.3\% to 52.9\% as $\rho$ increases from 10 to
70, while flagging rises from 21.0\% to 26.8\% and misreading from 5.4\% to
7.7\%\DIFdelbegin \DIFdel{. The }\DIFdelend \DIFaddbegin \DIFaddB{; the }\DIFaddend same overall pattern appears on MATH-500 and ARC. Misreading is also
\DIFdelbegin \DIFdel{closely }\DIFdelend related to failure: \DIFdelbegin \DIFdel{it occurs }\DIFdelend \DIFaddbegin \DIFaddB{in every configuration it is }\DIFaddend 1.2 to 2.0 times more \DIFdelbegin \DIFdel{often in incorrect answers than in correct ones }\DIFdelend \DIFaddbegin \DIFaddG{frequent
in GSM8K traces that end in a wrong answer, and 1.2 to 2.4 times on ARC; on
MATH-500, where misreads are rare, the relation is not consistent
}\DIFaddend (Table~\ref{tab:repair}). \DIFdelbegin %DIFDELCMD < 

%DIFDELCMD < %%%
\DIFdel{The proportion of typos classified as not used remains relatively }\DIFdelend \DIFaddbegin \DIFaddB{The share of typos that are not used stays }\DIFaddend stable
within each dataset \DIFdelbegin \DIFdel{, }\DIFdelend but differs across tasks\DIFdelbegin \DIFdel{: }\DIFdelend \DIFaddbegin \DIFaddB{, }\DIFaddend roughly 12\% on GSM8K, 13\% on
ARC, and 19\% on MATH-500\DIFdelbegin \DIFdel{. One possible explanation is that }\DIFdelend \DIFaddbegin \DIFaddB{, possibly because }\DIFaddend MATH-500 reasoning \DIFdelbegin \DIFdel{relies
more heavily on symbolic
manipulation, so the model }\DIFdelend \DIFaddbegin \DIFaddB{is more symbolic
and }\DIFaddend refers back to \DIFdelbegin \DIFdel{individual
words from the original }\DIFdelend \DIFaddbegin \DIFaddB{the wording of the }\DIFaddend question less often\DIFdelbegin \DIFdel{(Figure~\ref{fig:repair})}\DIFdelend .

\DIFdelbegin \DIFdel{In the LLM-as-a-judge analysis, a separate language model examined each
reasoning trace and assigned a score for repair behaviour. The judgewas shown
the corrupted words, the question, and the model's reasoning trace. It then gave
a score from 0 to 5: 0 means that the
model fully understood the meaning of the
question, while 5 means that it completely failed to understand the intended
meaning. We then calculated the mean score for different typo rates, real-word
ratios, and final-answer outcomes.
}%DIFDELCMD < 

%DIFDELCMD < %%%
\DIFdel{The results show that more typos generally make the question harder to
understand and repair. The average repair-understanding score increases as the
typo raterises, meaning that the model loses more of the original meaning when
more words are corrupted. Wrong answers also receive higher average
repair-understanding scores than
correct answers, supporting the connection
between misunderstanding }\DIFdelend \DIFaddbegin \DIFaddG{The judge's repair score (\S\ref{sec:setup}; Table~\ref{tab:judge}) tells the
same story. It rises with the typo rate, and wrong answers score higher than
correct ones (2.30 vs.\ 0.87 on GSM8K), linking misunderstanding of }\DIFaddend the
corrupted question \DIFdelbegin \DIFdel{and failure. The }\DIFdelend \DIFaddbegin \DIFaddB{to failure. It also rises with the }\DIFaddend real-word ratio \DIFdelbegin \DIFdel{has an additional effect: }\DIFdelend \DIFaddbegin \DIFaddB{in every
dataset and typo-rate cell, consistent with }\DIFaddend real-word \DIFdelbegin \DIFdel{typos look valid and can silently
change the meaning of a sentence, making them harder to
detect and repairthan
obviously invalid non-word typos (Table~\ref{tab:judge})}\DIFdelend \DIFaddbegin \DIFaddB{typos being harder to
repair, although we did not test this trend for significance}\DIFaddend .

\subsection{Self-doubt}
\DIFaddbegin \label{sec:doubt}
\DIFaddend 


\DIFaddbegin \begin{table}[t]
\centering
\footnotesize
\resizebox{\columnwidth}{!}{\begin{tabular}{l rr rr rr}
\toprule
& \multicolumn{2}{c}{\textbf{GSM8K}} & \multicolumn{2}{c}{\textbf{MATH-500}} & \multicolumn{2}{c}{\textbf{ARC-C}} \\
\cmidrule(lr){2-3}\cmidrule(lr){4-5}\cmidrule(lr){6-7}
\textbf{condition} & repair & doubt & repair & doubt & repair & doubt \\
\midrule
\multicolumn{7}{l}{\emph{clean}} \\
\quad --- & 0.13 & 4.50 & 0.03 & 3.66 & 0.03 & 4.69 \\
\midrule \multicolumn{7}{l}{\emph{by typo rate }$r$} \\
\quad 25\% & 0.88 & 4.86 & 0.35 & 3.90 & 0.97 & 5.12 \\
\quad 50\% & 1.27 & 5.12 & 0.35 & 3.98 & 1.03 & 5.31 \\
\quad 75\% & 1.74 & 5.64 & 0.48 & 4.01 & 1.35 & 5.63 \\
\midrule \multicolumn{7}{l}{\emph{by real-word ratio }$\rho$} \\
\quad 10\% & 1.14 & 5.21 & 0.32 & 3.87 & 0.94 & 5.33 \\
\quad 40\% & 1.32 & 5.29 & 0.39 & 4.03 & 1.10 & 5.27 \\
\quad 70\% & 1.42 & 5.12 & 0.47 & 3.99 & 1.31 & 5.45 \\
\midrule \multicolumn{7}{l}{\emph{by final answer}} \\
\quad correct & 0.87 & 4.69 & 0.27 & 3.75 & 0.88 & 4.94 \\
\quad wrong & 2.30 & 6.75 & 1.30 & 5.81 & 1.62 & 6.88 \\
\bottomrule
\end{tabular}}
\caption{Mean LLM-judge scores \DIFaddbegin \DIFaddB{over answered traces}\DIFaddend . \emph{repair} is the repair\DIFdelbegin \DIFdel{-understanding}\DIFdelend  score (0--5, higher means less of the intended meaning recovered); \emph{doubt}
is the self-doubt score (0--10, higher means more hesitation and looping)\DIFdelbegin \DIFdel{.
All three datasets were judged on the same main
generations reported everywhere else in this paper}\DIFdelend .}
\label{tab:judge}
\end{table}


\DIFaddend We measure self-doubt \DIFdelbegin \DIFdel{using }\DIFdelend \DIFaddbegin \DIFaddB{with }\DIFaddend two curated sets of expressions\DIFdelbegin \DIFdel{found in the
reasoning traces}\DIFdelend . The uncertainty set
\DIFdelbegin \DIFdel{includes hedging expressions }\DIFdelend \DIFaddbegin \DIFaddB{contains hedges }\DIFaddend such as \emph{maybe}, \emph{perhaps}, \emph{possibly},
\emph{might be}, and \emph{not sure}\DIFdelbegin \DIFdel{. The }\DIFdelend \DIFaddbegin \DIFaddB{; the }\DIFaddend second-guessing set captures explicit
reconsideration\DIFdelbegin \DIFdel{or
self-correction, with expressions }\DIFdelend \DIFaddbegin \DIFaddB{, }\DIFaddend such as \emph{wait}, \emph{actually}, \emph{reconsider},
\emph{let me recheck}, and \emph{double-check}. \DIFdelbegin \DIFdel{To account
for differences in trace length, we report the number of }\DIFdelend \DIFaddbegin \DIFaddB{We report }\DIFaddend matches per 1,000
reasoning words.

Uncertainty increases as the input becomes more corrupted. On GSM8K, \DIFdelbegin \DIFdel{for
example, uncertainty }\DIFdelend \DIFaddbegin \DIFaddB{it }\DIFaddend rises
from 7.4 markers per 1,000 words on clean input to 12.3 at a 25\% typo rate and
18.7 at \DIFdelbegin \DIFdel{a }\DIFdelend 75\%\DIFdelbegin \DIFdel{typo rate, when }\DIFdelend \DIFaddbegin \DIFaddB{, with }\DIFaddend the real-word ratio \DIFdelbegin \DIFdel{is }\DIFdelend held at 10\%. The same \DIFdelbegin \DIFdel{increasing pattern appears on
ARC and }\DIFdelend \DIFaddbegin \DIFaddB{trend holds on
}\DIFaddend MATH-500 \DIFdelbegin \DIFdel{.
Increasing the real-word ratio at }\DIFdelend \DIFaddbegin \DIFaddB{and ARC at the higher typo rates, although on ARC the mildest rate is
not above clean. At }\DIFaddend a fixed typo rate\DIFdelbegin \DIFdel{leads to an additional rise
in }\DIFdelend \DIFaddbegin \DIFaddB{, a higher real-word ratio tends to add
further }\DIFaddend uncertainty (Figure~\ref{fig:doubt}; Table~\ref{tab:support}).

Second-guessing \DIFdelbegin \DIFdel{shows a different pattern}\DIFdelend \DIFaddbegin \DIFaddB{behaves differently}\DIFaddend . Its density \DIFdelbegin \DIFdel{increases only slightly between clean and corrupted inputs and then remains relatively stable }\DIFdelend \DIFaddbegin \DIFaddG{rises only slightly on GSM8K,
not at all on MATH-500 and ARC, and stays roughly flat }\DIFaddend across typo rates and
real-word ratios. \DIFdelbegin \DIFdel{Although the total }\DIFdelend \DIFaddbegin \DIFaddB{The }\DIFaddend number of second-guessing expressions per trace \DIFdelbegin \DIFdel{increases,
this is largely explained by the fact that
}\DIFdelend \DIFaddbegin \DIFaddB{does grow,
but mostly because }\DIFaddend corrupted inputs produce longer \DIFdelbegin \DIFdel{reasoning traces. In other words, typos }\DIFdelend \DIFaddbegin \DIFaddB{traces. Typos thus }\DIFaddend make the
model express more uncertainty \DIFdelbegin \DIFdel{, but do not appear to make }\DIFdelend \DIFaddbegin \DIFaddB{without making }\DIFaddend it reconsider its reasoning more
often\DIFdelbegin \DIFdel{(Figure~\ref{fig:doubt})}\DIFdelend .

The \DIFdelbegin \DIFdel{difference becomes }\DIFdelend \DIFaddbegin \DIFaddB{contrast is }\DIFaddend clearer when we separate correct and incorrect answers. For the
same \DIFdelbegin \DIFdel{typo rate and real-word ratio, uncertainty expressions are
}\DIFdelend \DIFaddbegin \DIFaddB{configuration, uncertainty is }\DIFaddend consistently more frequent in traces that end
with a wrong answer\DIFdelbegin \DIFdel{. On GSM8K,
uncertainty density is about 1.4 }\DIFdelend \DIFaddbegin \DIFaddG{: 1.6 }\DIFaddend to 2.1 times \DIFdelbegin \DIFdel{higher in incorrect traces}\DIFdelend \DIFaddbegin \DIFaddB{on GSM8K}\DIFaddend , with similar patterns on ARC and
MATH-500\DIFdelbegin \DIFdel{. Second-guessing shows a much weaker
difference between correct and incorrect answers }\DIFdelend \DIFaddbegin \DIFaddB{, while second-guessing differs much less }\DIFaddend (Table~\ref{tab:support}).
\DIFdelbegin %DIFDELCMD < 

%DIFDELCMD < %%%
\DIFdel{Overall, uncertainty language is strongly associated with model failure, while
explicit self-correction is not . The model expresses more doubt as the input
becomes noisier and also tends to express more doubt when it is about to give a
wrong answer. However, this relationship is not unique to typo-corrupted inputs}\DIFdelend \DIFaddbegin \DIFaddB{This relation is not unique to typos}\DIFaddend : even on clean GSM8K questions, incorrect
answers \DIFdelbegin \DIFdel{already }\DIFdelend contain substantially more uncertainty language than correct ones.

\DIFdelbegin \DIFdel{In the LLM-as-a-judge analysis, a separate language model examined each
reasoning trace and assigned a }\DIFdelend \DIFaddbegin \DIFaddB{The judge's }\DIFaddend self-doubt \DIFdelbegin \DIFdel{score from 0 to 10. A score of 0 means
that the model followed a clear, straight-line solution without hesitation,
while 10 means that the trace contained pervasive uncertainty, repeated
second-guessing, reversals, or looping. The judgewas shown the corrupted words,
the question, and the model's reasoning trace (like in 5.3) .
}%DIFDELCMD < 

%DIFDELCMD < %%%
\DIFdel{On }\DIFdelend \DIFaddbegin \DIFaddB{score (Table~\ref{tab:judge}) rises with the typo rate on
}\DIFaddend all three datasets, \DIFdelbegin \DIFdel{the average self-doubt score increases as the typo rate
rises, indicating that more corrupted words cause more hesitation and re-evaluation. Typo traces also }\DIFdelend \DIFaddbegin \DIFaddB{and typo traces }\DIFaddend score higher than clean \DIFdelbegin \DIFdel{traces }\DIFdelend \DIFaddbegin \DIFaddB{ones }\DIFaddend everywhere. The
real-word ratio, \DIFdelbegin \DIFdel{by contrast, has almost no effect: at a fixed typo rate the score moves by at most about 0.1 point in either direction. The amount of
corruption changes how much the model hesitates, while }\DIFdelend \DIFaddbegin \DIFaddG{however, shows no consistent direction: pooled over typo rates,
the score varies by less than 0.2 points on each dataset. The judge therefore
confirms that more corruption brings more hesitation but, unlike the marker
counts, does not register an additional effect of }\DIFaddend the kind of corruption\DIFdelbegin \DIFdel{does not (Table~\ref{tab:judge})}\DIFdelend .


\DIFdelbegin %DIFDELCMD < \FLOATPLACEHOLDER{tab:judge}
%DIFDELCMD < 

%DIFDELCMD < %%%
\DIFdelend \DIFdelbegin %DIFDELCMD < \begin{figure*}[tbp]
%DIFDELCMD < %%%
\DIFdelendFL \DIFaddbeginFL \begin{figure*}[t]
\DIFaddendFL \centering
\includegraphics[width=\textwidth]{figures/fig_doubt.pdf}
\caption{\textbf{Self-doubt under typos.} Marker density per 1{,}000
reasoning words across the grid. Uncertainty markers (solid) \DIFdelbeginFL \DIFdelFL{climb steadily
}\DIFdelendFL \DIFaddbeginFL \DIFaddBFL{rise }\DIFaddendFL with
corruption on all three datasets, \DIFaddbeginFL \DIFaddBFL{most steeply on GSM8K, }\DIFaddendFL while second-guessing
markers (dashed) stay \DIFaddbeginFL \DIFaddBFL{roughly }\DIFaddendFL flat\DIFdelbeginFL \DIFdelFL{, so }\DIFdelendFL \DIFaddbeginFL \DIFaddBFL{: }\DIFaddendFL typos make the model express more doubt
without making it verify more. The same densities split by final answer are in
Table~\ref{tab:support}.}
\label{fig:doubt}
\end{figure*}


\subsection{\DIFdelbegin \DIFdel{Every }\DIFdelend \DIFaddbegin \DIFaddB{No }\DIFaddend lightweight mitigation \DIFdelbegin \DIFdel{fails}\DIFdelend \DIFaddbegin \DIFaddB{reliably helps}\DIFaddend }
\DIFaddbegin \label{sec:fixes}
\DIFaddend 

We tested the three mitigation strategies \DIFdelbegin \DIFdel{proposed in our plan , all }\DIFdelend \DIFaddbegin \DIFaddB{from our plan }\DIFaddend on GSM8K
(\DIFdelbegin \DIFdel{Figure~\ref{fig:fixes}, }\DIFdelend Table~\ref{tab:fixes}).

\paragraph{Warn.}
\DIFdelbegin \DIFdel{We first added a short warning that the input }\DIFdelend \DIFaddbegin \DIFaddB{Prefixing the prompt with ``Note: the text }\DIFaddend may contain typos.\DIFdelbegin \DIFdel{Across the nine
typo configurations, this }\DIFdelend \DIFaddbegin \DIFaddB{'' }\DIFaddend changed
accuracy by only $-0.5$ percentage points on average \DIFaddbegin \DIFaddB{across the nine typo
configurations}\DIFaddend , with effects ranging from $-4.3$ to $+3.5$ points. \DIFdelbegin \DIFdel{In other words,
simply warning the model does not improve performance. }\DIFdelend This is
consistent with \DIFdelbegin \DIFdel{the analysis in Section~5.3}\DIFdelend \DIFaddbegin \DIFaddB{\S\ref{sec:repair}}\DIFaddend , where the model often already showed signs
of noticing that the input was corrupted\DIFdelbegin \DIFdel{. The main limitation therefore }\DIFdelend \DIFaddbegin \DIFaddB{: the limitation }\DIFaddend does not appear to be
awareness of the typo itself.

\paragraph{Rewrite.}
\DIFdelbegin \DIFdel{We next instructed }\DIFdelend \DIFaddbegin \DIFaddB{Instructing }\DIFaddend the model to first rewrite the question with the typos corrected
and only then solve it \DIFdelbegin \DIFdel{. This }\DIFdelend performed worse than the original \DIFdelbegin \DIFdel{setup
}\DIFdelend \DIFaddbegin \DIFaddB{prompt }\DIFaddend in all nine
configurations, reducing accuracy by 4.7 percentage points on average and by \DIFdelbegin \DIFdel{as much as }\DIFdelend \DIFaddbegin \DIFaddB{up
to }\DIFaddend 7.5 points. Median reasoning length also \DIFdelbegin \DIFdel{dropped from }\DIFdelend \DIFaddbegin \DIFaddB{fell, from about }\DIFaddend 1,\DIFdelbegin \DIFdel{214 to about 730 tokens }\DIFdelend \DIFaddbegin \DIFaddG{300 to about 740
tokens (median over the nine typo configurations)}\DIFaddend . One possible explanation is
that the rewrite step forces the model to commit to a single interpretation of
the corrupted question too early\DIFdelbegin \DIFdel{. If }\DIFdelend \DIFaddbegin \DIFaddB{; if }\DIFaddend the reconstruction is wrong, the model
\DIFdelbegin \DIFdel{then }\DIFdelend reasons about the wrong \DIFdelbegin \DIFdel{version of the }\DIFdelend problem with little opportunity to recover.
\DIFdelbegin \DIFdel{The additional
rewrite step may also shift part of the generation budget away from the
reasoning process itself.
}\DIFdelend 

\paragraph{Spell check.}
We applied an external \DIFdelbegin \DIFdel{spellchecker }\DIFdelend \DIFaddbegin \DIFaddB{spell checker (}\texttt{\DIFaddB{pyspellchecker}}\DIFaddB{, edit distance 2,
prose words only) }\DIFaddend to the same \DIFdelbegin \DIFdel{500 GSM8K }\DIFdelend questions before sending them to the model. A
typo \DIFdelbegin \DIFdel{was considered }\DIFdelend \DIFaddbegin \DIFaddB{counts as }\DIFaddend recovered only if it \DIFdelbegin \DIFdel{was
}\DIFdelend \DIFaddbegin \DIFaddB{is }\DIFaddend restored exactly to the original word.
Recovery dropped as the share of real-word typos increased: from 54.9\% at
$\rho{=}10$ to 39.9\% at $\rho{=}40$ and 25.8\% at $\rho{=}70$. \DIFdelbegin \DIFdel{However, this }\DIFdelend \DIFaddbegin \DIFaddB{This }\DIFaddend lexical
recovery did not translate into a consistent accuracy gain\DIFdelbegin \DIFdel{. Across the nine typo
configurations, the spellchecker }\DIFdelend \DIFaddbegin \DIFaddB{: the spell checker
}\DIFaddend changed accuracy by \DIFdelbegin \DIFdel{only }\DIFdelend $-0.1$ percentage points on average, improving five
conditions and hurting four\DIFdelbegin \DIFdel{; accuracy }\DIFdelend \DIFaddbegin \DIFaddG{, with 62.2\% }\DIFaddend in the hardest condition\DIFdelbegin \DIFdel{was 62.2\% (Table~\ref{tab:fixes}). }%DIFDELCMD < 

%DIFDELCMD < %%%
\DIFdel{This is likely because spellcheckers are much better at detecting }\DIFdelend \DIFaddbegin \DIFaddB{. Spell
checkers detect }\DIFaddend non-word errors \DIFaddbegin \DIFaddB{far better }\DIFaddend than real-word substitutions, \DIFdelbegin \DIFdel{which already look like valid words. In
some cases, they }\DIFdelend \DIFaddbegin \DIFaddB{and
}\DIFaddend may even replace a non-word typo with \DIFdelbegin \DIFdel{the }\DIFdelend \DIFaddbegin \DIFaddB{a }\DIFaddend wrong valid word, \DIFdelbegin \DIFdel{effectively creating the type }\DIFdelend \DIFaddbegin \DIFaddB{creating the kind }\DIFaddend of
error that \DIFdelbegin \DIFdel{Section~5.2 }\DIFdelend \DIFaddbegin \DIFaddB{\S\ref{sec:realword} }\DIFaddend found to be more harmful.


\DIFdelbegin \DIFdel{Taken together, these results suggest that the main difficulty is not simply
detecting that the input contains errors. None of the three strategies reliably
recovers the information that was lost during corruption: the identity of the
original word. Warnings provide little additional information, an explicit
rewrite step can commit the model to an incorrect reconstruction, and
dictionary-based correction is least effective precisely when the corruption
results in another valid word.
}%DIFDELCMD < 

%DIFDELCMD < \FLOATPLACEHOLDER{fig:fixes}
\begin{center}\fbox{\parbox{0.9\linewidth}{\small\DIFdel{Removed figure: \textbf{Prompt-level mitigations on GSM8K.} \emph{(a)} answered-only
accuracy for the unmodified prompt, the typo warning and the rewrite-first
instruction. \emph{(b)} each fix as a change from the unmodified prompt. Warning
the model is null on average and the rewrite instruction is consistently
harmful.}}}\end{center}
%DIFDELCMD < 

%DIFDELCMD < %%%
\DIFdelend \DIFdelbegin %DIFDELCMD < \begin{table}[tbp]
%DIFDELCMD < %%%
\DIFdelendFL \DIFaddbeginFL \begin{table}[t]
\DIFaddendFL \centering
\footnotesize
\DIFaddbeginFL \setlength{\tabcolsep}{4pt}
\DIFaddendFL \adjustbox{max width=\linewidth}{\begin{tabular}{lr rr rr rr}
\toprule
& \textbf{no fix} & \multicolumn{2}{c}{\textbf{warn}} & \multicolumn{2}{c}{\textbf{rewrite}} & \multicolumn{2}{c}{\textbf{spell check}} \\
\cmidrule(lr){3-4}\cmidrule(lr){5-6}\cmidrule(lr){7-8}
\textbf{$r$/$\rho$} & acc & acc & $\Delta$ & acc & $\Delta$ & acc & $\Delta$ \\
\midrule
25/10 & 85.5 & 85.7 & +0.2 & 83.8 & -1.7 & 87.0 & +1.5 \\
25/40 & 87.0 & 87.2 & +0.2 & 82.9 & -4.1 & 86.1 & -0.9 \\
25/70 & 86.7 & 85.5 & -1.2 & 80.2 & -6.6 & 85.3 & -1.4 \\
50/10 & 79.6 & 79.4 & -0.2 & 76.3 & -3.3 & 79.0 & -0.6 \\
50/40 & 79.0 & 74.7 & -4.3 & 71.5 & -7.5 & 75.0 & -4.0 \\
50/70 & 72.8 & 74.0 & +1.2 & 67.9 & -4.9 & 73.2 & +0.4 \\
75/10 & 73.8 & 70.6 & -3.2 & 66.9 & -6.9 & 75.3 & +1.5 \\
75/40 & 67.3 & 66.3 & -0.9 & 63.8 & -3.5 & 68.1 & +0.9 \\
75/70 & 60.4 & 63.9 & +3.5 & 56.9 & -3.5 & 62.2 & +1.8 \\
\midrule \textbf{mean} & & & \textbf{-0.5} & & \textbf{-4.7} & & \textbf{-0.1} \\
\bottomrule
\end{tabular}}
\caption{GSM8K \DIFaddbeginFL \DIFaddBFL{answered-only }\DIFaddendFL accuracy (\%) under each \DIFdelbeginFL \DIFdelFL{prompt-level }\DIFdelendFL mitigation, per
configuration, with the change against the unmodified prompt \DIFaddbeginFL \DIFaddBFL{(}\emph{\DIFaddBFL{no fix}}\DIFaddBFL{)}\DIFaddendFL .
The mean row is taken over the nine typo configurations.}
\label{tab:fixes}
\end{table}


% =====================================================================
\section{Discussion}
%DIF <  AUTHOR TEXT - verbatim.
% =====================================================================

The most consistent finding across our analyses is that real-word typos are more
harmful than non-word typos\DIFdelbegin \DIFdel{. This pattern appears }\DIFdelend \DIFaddbegin \DIFaddB{, }\DIFaddend in the per-typo regression, the paired comparisons,
and the \DIFdelbegin \DIFdel{spellchecker results. At first, this may seem
surprising, }\DIFdelend \DIFaddbegin \DIFaddB{spell-checker results. One might expect the opposite, }\DIFaddend since non-word
typos \DIFdelbegin \DIFdel{often produce unusual tokenizationpatterns and
might therefore be expected to be more disruptive}\DIFdelend \DIFaddbegin \DIFaddB{disrupt tokenization}\DIFaddend . Our results suggest \DIFdelbegin \DIFdel{the
opposite. A }\DIFdelend \DIFaddbegin \DIFaddB{that a }\DIFaddend non-word typo is visibly
unusual and \DIFdelbegin \DIFdel{can serve as a signal }\DIFdelend \DIFaddbegin \DIFaddB{signals }\DIFaddend that something is wrong, \DIFdelbegin \DIFdel{while }\DIFdelend \DIFaddbegin \DIFaddB{whereas }\DIFaddend a real-word typo is fluent,
tokenized normally, and may fit the sentence well enough \DIFdelbegin \DIFdel{to go unnoticed even when its meaning is
incorrect. }%DIFDELCMD < 

%DIFDELCMD < %%%
\DIFdel{This also helps explain the limited success of }\DIFdelend \DIFaddbegin \DIFaddB{that its wrong meaning
is taken at face value. The same property makes it hard for }\DIFaddend dictionary-based
correction\DIFdelbegin \DIFdel{. Such
methods are effective mainly when the corrupted form is not a valid word.
Real-word substitutions, however, are already present in the dictionary and are
therefore much harder to detect. In our spellchecker experiment, recovery
dropped substantially as the proportion of real-word typos increased. In some
cases, the spellchecker also replaced a non-word typo with an incorrect valid
word, effectively turning a visible error into the type of corruption we found
to be more damaging.
}\DIFdelend \DIFaddbegin \DIFaddB{, which can only act on forms that are not valid words.
}\DIFaddend 

The trace analysis points to a related conclusion. The model often \DIFdelbegin \DIFdel{appears to
notice }\DIFdelend \DIFaddbegin \DIFaddB{notices }\DIFaddend that
the input is unusual: uncertainty language increases, \DIFdelbegin \DIFdel{suspicious
}\DIFdelend \DIFaddbegin \DIFaddB{corrupted }\DIFaddend words are
flagged more often, and \DIFdelbegin \DIFdel{reasoning }\DIFdelend traces become longer\DIFdelbegin \DIFdel{as corruption
increases. However, this }\DIFdelend \DIFaddbegin \DIFaddB{. This }\DIFaddend additional effort does not
reliably lead to recovery\DIFdelbegin \DIFdel{.
The model expresses more uncertainty, but does not appear to verify its
interpretation more effectively.
}%DIFDELCMD < 

%DIFDELCMD < %%%
\DIFdel{This helps explain why the simple mitigation strategies we tested }\DIFdelend \DIFaddbegin \DIFaddB{, which also explains why the mitigations }\DIFaddend had little
effect\DIFdelbegin \DIFdel{. Warning the model adds information it }\DIFdelend \DIFaddbegin \DIFaddB{: a warning adds information the model }\DIFaddend often already has, while \DIFdelbegin \DIFdel{forcing
it to rewrite the question can make it commit too }\DIFdelend \DIFaddbegin \DIFaddB{a forced
rewrite commits it }\DIFaddend early to a single\DIFdelbegin \DIFdel{interpretation. If that reconstruction is wrong, the subsequent reasoning is
based on the wrongversion of the problem.
}%DIFDELCMD < 

%DIFDELCMD < %%%
\DIFdel{Taken together, these results suggest that the }\DIFdelend \DIFaddbegin \DIFaddB{, possibly wrong, interpretation. The }\DIFaddend main
difficulty is \DIFaddbegin \DIFaddB{therefore }\DIFaddend not detecting that a typo exists, but recovering the
intended word when the corruption is ambiguous. A more promising direction may
\DIFdelbegin \DIFdel{therefore be to preserve multiple possible interpretations }\DIFdelend \DIFaddbegin \DIFaddB{be to keep multiple interpretations open }\DIFaddend during reasoning and resolve \DIFdelbegin \DIFdel{the ambiguity }\DIFdelend \DIFaddbegin \DIFaddB{them
}\DIFaddend later using the structure of the problem\DIFdelbegin \DIFdel{, rather than forcing an early correction. }%DIFDELCMD < 

%DIFDELCMD < %%%
\DIFdelend \DIFaddbegin \DIFaddB{. }\DIFaddend Finally, longer \DIFdelbegin \DIFdel{reasoning }\DIFdelend traces should not be
\DIFdelbegin \DIFdel{interpreted }\DIFdelend \DIFaddbegin \DIFaddB{read }\DIFaddend as evidence of successful self-correction\DIFdelbegin \DIFdel{. In our results, incorrect answers tend to }\DIFdelend \DIFaddbegin \DIFaddB{: incorrect answers }\DIFaddend produce
longer traces even on clean input.
\DIFdelbegin \DIFdel{This suggests that additional reasoning often
reflects difficulty rather than successful recovery.
}\DIFdelend 

% =====================================================================
\section{Limitations}
%DIF <  AUTHOR TEXT - verbatim.
% =====================================================================

\DIFdelbegin \paragraph{\DIFdel{One model, one sample.}}
%DIFAUXCMD
\addtocounter{paragraph}{-1}%DIFAUXCMD
\DIFdelend \DIFaddbegin \textbf{\DIFaddB{One model, one sample.}}
\DIFaddend All headline results come from a single 7B reasoning model at one decoding
setting with $n{=}1$ sample per question. Our plan called for several samples per
question, but the API budget did not permit it, so we cannot separate typo
effects from decoding variance at the level of an individual question. \DIFdelbegin \DIFdel{Still,
the }\DIFdelend \DIFaddbegin \DIFaddB{The
}\DIFaddend overall effects are much larger than expected sampling noise (McNemar
$p < 10^{-6}$ on $n{=}480$ paired items), but \DIFdelbegin \DIFdel{generalisation }\DIFdelend \DIFaddbegin \DIFaddB{generalization }\DIFaddend to larger reasoning
models is untested.

\DIFdelbegin \paragraph{\DIFdel{Dataset substitution.}}
%DIFAUXCMD
\addtocounter{paragraph}{-1}%DIFAUXCMD
\DIFdelend \DIFaddbegin \textbf{\DIFaddB{Dataset substitution.}}
\DIFaddend We originally planned to use GPQA-Diamond, but replaced it with ARC-Challenge,
which is also a multiple-choice science dataset. On clean GPQA-Diamond, the
model reached only about 52\% accuracy. Since random guessing gives 25\%, this
leaves limited room to measure further drops under corruption: as performance
gets closer to chance, changes in accuracy become harder to interpret as
reasoning failures rather than guessing. ARC-Challenge is easier and its answer
options may help the model recover some information, so our multiple-choice
results may underestimate the effect of typos on harder science questions.

\DIFdelbegin \paragraph{\DIFdel{Mitigations and human validation.}}
%DIFAUXCMD
\addtocounter{paragraph}{-1}%DIFAUXCMD
\DIFdelend \DIFaddbegin \textbf{\DIFaddB{Scope of the mitigation study.}}
\DIFaddend We tested the three mitigation methods only on GSM8K because the experiments
required additional paid API calls and our budget was limited. We chose \DIFdelbegin \DIFdel{to focus
on }\DIFdelend GSM8K
because it showed the clearest drop in answer quality under typo corruption,
making it the most informative dataset for evaluating whether these mitigation
strategies could recover performance.

% =====================================================================
\section{AI Disclosure and Reflection}
%DIF <  AUTHOR TEXT - verbatim.
% =====================================================================

\paragraph{What we used AI for.}
We used GitHub Copilot with Claude primarily to support coding tasks during the
project. It helped us build the typo-generation and inference pipelines, write
analysis scripts, generate matplotlib code for figures, and improve the wording
of some parts of the paper. All methodological and experimental decisions,
including the two-axis corruption grid, typo-type weights, answer-extraction
rules, and judge rubric, were made by us. The analysis scripts were written with
Claude's assistance based on these decisions, and all reported results were
generated from the raw generation logs using those scripts.

\paragraph{Reflection.}
AI tools were especially useful for speeding up repetitive coding and analysis
tasks and for keeping the implementation consistent across experiments. However,
generated code and outputs still required careful checking, particularly for
numerical results. We therefore verified reported results against the raw data
and used a consistent scoring procedure throughout the analysis. Overall, the
project highlighted the importance of reproducibility and independent
verification when AI tools are used in research.

% =====================================================================
\section{Conclusion}
%DIF <  AUTHOR TEXT - verbatim.
% =====================================================================

Typos have a clear and consistent effect on reasoning performance. On GSM8K,
accuracy dropped by as much as 31 percentage points, with larger declines as the
corruption rate increased\DIFdelbegin \DIFdel{. The }\DIFdelend \DIFaddbegin \DIFaddB{, and the }\DIFaddend strongest effect came from typos that turned
one valid English word into another. \DIFdelbegin %DIFDELCMD < 

%DIFDELCMD < %%%
\DIFdel{Our analysis suggests that the }\DIFdelend \DIFaddbegin \DIFaddB{The }\DIFaddend main problem is not that the model
fails to notice corrupted input\DIFdelbegin \DIFdel{. In many cases, it appears to recognize }\DIFdelend \DIFaddbegin \DIFaddB{: it often recognizes }\DIFaddend that something is wrong,
but \DIFdelbegin \DIFdel{it still }\DIFdelend cannot reliably recover the intended meaning. \DIFdelbegin \DIFdel{The reasoning
}\DIFdelend \DIFaddbegin \DIFaddB{Its }\DIFaddend traces show more
uncertainty \DIFdelbegin \DIFdel{than actual }\DIFdelend \DIFaddbegin \DIFaddB{rather than more }\DIFaddend verification, and \DIFdelbegin \DIFdel{the mitigation
methods we tested, including prompting strategies and }\DIFdelend \DIFaddbegin \DIFaddB{neither prompting strategies nor
}\DIFaddend dictionary-based correction \DIFdelbegin \DIFdel{, did not consistently improve performance. }%DIFDELCMD < 

%DIFDELCMD < %%%
\DIFdel{Overall, these results suggest that robustness }\DIFdelend \DIFaddbegin \DIFaddB{consistently improved performance. Robustness }\DIFaddend to
everyday typing errors should \DIFaddbegin \DIFaddB{therefore }\DIFaddend be treated as a \DIFdelbegin \DIFdel{separate capability that needs }\DIFdelend \DIFaddbegin \DIFaddB{capability }\DIFaddend to be
developed explicitly, rather than \DIFdelbegin \DIFdel{something that can be }\DIFdelend assumed from strong benchmark performance.

All code is available in the project repository, and the raw generation logs are
hosted on the Hugging Face Hub (\texttt{ronshtricker/\allowbreak typo-reasoning-results}).
All reported results can be reproduced by running \texttt{download\_data.py}
followed by \texttt{run\_all.py}.

%DIF <  =====================================================================
\DIFdelbegin %DIFDELCMD < 

%DIFDELCMD < %%%
\DIFdelend \begin{thebibliography}{17}
\providecommand{\natexlab}[1]{#1}

\bibitem[{Belinkov and Bisk(2018)}]{belinkov2018}
Yonatan Belinkov and Yonatan Bisk. 2018.
\newblock Synthetic and natural noise both break neural machine translation.
\newblock In \emph{International Conference on Learning Representations}.

\bibitem[{Chai et~al.(2024)Chai, Fang, Peng, and Li}]{chai2024}
Yekun Chai, Yewei Fang, Qiwei Peng, and Xuhong Li. 2024.
\newblock Tokenization falling short: On subword robustness in large language
  models.
\newblock In \emph{Findings of the Association for Computational Linguistics:
  EMNLP 2024}.

\bibitem[{Clark et~al.(2018)Clark, Cowhey, Etzioni, Khot, Sabharwal, Schoenick,
  and Tafjord}]{clark2018}
Peter Clark, Isaac Cowhey, Oren Etzioni, Tushar Khot, Ashish Sabharwal, Carissa
  Schoenick, and Oyvind Tafjord. 2018.
\newblock Think you have solved question answering? try {ARC}, the {AI2}
  reasoning challenge.
\newblock \emph{arXiv preprint arXiv:1803.05457}.

\bibitem[{Cobbe et~al.(2021)Cobbe, Kosaraju, Bavarian, Chen, Jun, Kaiser
  et~al.}]{cobbe2021}
Karl Cobbe, Vineet Kosaraju, Mohammad Bavarian, Mark Chen, Heewoo Jun, Lukasz
  Kaiser, et~al. 2021.
\newblock Training verifiers to solve math word problems.
\newblock \emph{arXiv preprint arXiv:2110.14168}.

\bibitem[{Ebrahimi et~al.(2018)Ebrahimi, Rao, Lowd, and Dou}]{ebrahimi2018}
Javid Ebrahimi, Anyi Rao, Daniel Lowd, and Dejing Dou. 2018.
\newblock {HotFlip}: White-box adversarial examples for text classification.
\newblock In \emph{Proceedings of the 56th Annual Meeting of the Association
  for Computational Linguistics}.

\bibitem[{Gan et~al.(2024)Gan, Zhao, Cheng, Mao, Goyal, Kawaguchi, Kan, and
  Shieh}]{gan2024}
Esther Gan, Yiran Zhao, Liying Cheng, Yancan Mao, Anirudh Goyal, Kenji
  Kawaguchi, Min-Yen Kan, and Michael Shieh. 2024.
\newblock Reasoning robustness of {LLM}s to adversarial typographical errors.
\newblock In \emph{Proceedings of the 2024 Conference on Empirical Methods in
  Natural Language Processing}.

\bibitem[{Grattafiori et~al.(2024)Grattafiori, Dubey, Jauhri, Pandey
  et~al.}]{llama3}
Aaron Grattafiori, Abhimanyu Dubey, Abhinav Jauhri, Abhinav Pandey, et~al.
  2024.
\newblock The {Llama} 3 herd of models.
\newblock \emph{arXiv preprint arXiv:2407.21783}.

\bibitem[{Guo et~al.(2025)Guo, Yang, Zhang, Song, Zhang, Xu
  et~al.}]{deepseekr1}
Daya Guo, Dejian Yang, Haowei Zhang, Junxiao Song, Ruoyu Zhang, Runxin Xu,
  et~al. 2025.
\newblock {DeepSeek-R1}: Incentivizing reasoning capability in {LLM}s via
  reinforcement learning.
\newblock \emph{arXiv preprint arXiv:2501.12948}.

\bibitem[{Hendrycks et~al.(2021)Hendrycks, Burns, Kadavath, Arora, Basart,
  Tang, Song, and Steinhardt}]{hendrycks2021}
Dan Hendrycks, Collin Burns, Saurav Kadavath, Akul Arora, Steven Basart, Eric
  Tang, Dawn Song, and Jacob Steinhardt. 2021.
\newblock Measuring mathematical problem solving with the {MATH} dataset.
\newblock In \emph{Proceedings of the Neural Information Processing Systems
  Track on Datasets and Benchmarks}.

\bibitem[{Huang et~al.(2024)Huang, Chen, Mishra, Zheng, Yu, Song, and
  Zhou}]{huang2024}
Jie Huang, Xinyun Chen, Swaroop Mishra, Huaixiu~Steven Zheng, Adams~Wei Yu,
  Xinying Song, and Denny Zhou. 2024.
\newblock Large language models cannot self-correct reasoning yet.
\newblock In \emph{The Twelfth International Conference on Learning
  Representations}.

\bibitem[{Jin et~al.(2020)Jin, Jin, Zhou, and Szolovits}]{jin2020}
Di~Jin, Zhijing Jin, Joey~Tianyi Zhou, and Peter Szolovits. 2020.
\newblock Is {BERT} really robust? a strong baseline for natural language
  attack on text classification and entailment.
\newblock In \emph{Proceedings of the AAAI Conference on Artificial
  Intelligence}.

\bibitem[{Lightman et~al.(2024)Lightman, Kosaraju, Burda, Edwards, Baker, Lee,
  Leike, Schulman, Sutskever, and Cobbe}]{lightman2024}
Hunter Lightman, Vineet Kosaraju, Yura Burda, Harri Edwards, Bowen Baker, Teddy
  Lee, Jan Leike, John Schulman, Ilya Sutskever, and Karl Cobbe. 2024.
\newblock Let's verify step by step.
\newblock In \emph{The Twelfth International Conference on Learning
  Representations}.

\bibitem[{Pruthi et~al.(2019)Pruthi, Dhingra, and Lipton}]{pruthi2019}
Danish Pruthi, Bhuwan Dhingra, and Zachary~C. Lipton. 2019.
\newblock Combating adversarial misspellings with robust word recognition.
\newblock In \emph{Proceedings of the 57th Annual Meeting of the Association
  for Computational Linguistics}.

\bibitem[{Turpin et~al.(2023)Turpin, Michael, Perez, and Bowman}]{turpin2023}
Miles Turpin, Julian Michael, Ethan Perez, and Samuel~R. Bowman. 2023.
\newblock Language models don't always say what they think: Unfaithful
  explanations in chain-of-thought prompting.
\newblock In \emph{Advances in Neural Information Processing Systems}.

\bibitem[{Wei et~al.(2022)Wei, Wang, Schuurmans, Bosma, Ichter, Xia, Chi, Le,
  and Zhou}]{wei2022}
Jason Wei, Xuezhi Wang, Dale Schuurmans, Maarten Bosma, Brian Ichter, Fei Xia,
  Ed~Chi, Quoc Le, and Denny Zhou. 2022.
\newblock Chain-of-thought prompting elicits reasoning in large language
  models.
\newblock In \emph{Advances in Neural Information Processing Systems}.

\bibitem[{Zhao et~al.(2026)Zhao, Liu, Altinger, Sch{\"u}tze, and
  Hedderich}]{zhao2026}
Raoyuan Zhao, Yihong Liu, Lena Altinger, Hinrich Sch{\"u}tze, and Michael~A.
  Hedderich. 2026.
\newblock Evaluating robustness of large language models against multilingual
  typographical errors.
\newblock In \emph{Proceedings of the 64th Annual Meeting of the Association
  for Computational Linguistics (Volume 1: Long Papers)}, pages 16059--16078,
  San Diego, California, United States.

\bibitem[{Zhu et~al.(2024)Zhu, Wang, Zhou, Wang, Chen, Wang, Yang, Ye, Zhang,
  Gong, and Xie}]{zhu2024}
Kaijie Zhu, Jindong Wang, Jiaheng Zhou, Zichen Wang, Hao Chen, Yidong Wang,
  Linyi Yang, Wei Ye, Yue Zhang, Neil Gong, and Xing Xie. 2024.
\newblock {PromptBench}: Towards evaluating the robustness of large language
  models on adversarial prompts.
\newblock In \emph{Proceedings of the 62nd Annual Meeting of the Association
  for Computational Linguistics}.

\end{thebibliography}


%DIF <  =====================================================================
\DIFdelbegin %DIFDELCMD < \appendix
%DIFDELCMD < %%%
\DIFdelend \clearpage
\onecolumn
\DIFdelbegin %DIFDELCMD < \raggedbottom
%DIFDELCMD < %%%
\section*{\DIFdel{Appendix}}
%DIFAUXCMD
%DIFDELCMD < \vspace{0.5em}
%DIFDELCMD < \renewcommand{\thetable}{A\arabic{table}}
%DIFDELCMD < \setcounter{table}{0}
%DIFDELCMD < \setlength{\tabcolsep}{2.6pt}
%DIFDELCMD < %%%
\DIFdelend %DIF > -------------------------------------------------------------------------------
\DIFaddbegin \appendix
\DIFaddend 

%DIF <  =====================================================================
\DIFaddbegin \section{\DIFaddB{Full accuracy grid}}
\label{app:grid}
\DIFaddend 

\DIFdelbegin %DIFDELCMD < \begin{table}[H]
%DIFDELCMD < %%%
\DIFdelendFL \DIFaddbeginFL \begin{table}[H]
\DIFaddendFL \centering\scriptsize\DIFaddbeginFL \setlength{\tabcolsep}{2.6pt}
\DIFaddendFL \adjustbox{max width=\linewidth}{\begin{tabular}{l rrrrrr rrrrrr rrrrrr}
\toprule
& \multicolumn{6}{c}{\textbf{GSM8K}} & \multicolumn{6}{c}{\textbf{MATH-500}} & \multicolumn{6}{c}{\textbf{ARC-Challenge}} \\
\cmidrule(lr){2-7}\cmidrule(lr){8-13}\cmidrule(lr){14-19}
\textbf{config} & acc & $n_a$ & tok & r$\to$w & w$\to$r & $p$ & acc & $n_a$ & tok & r$\to$w & w$\to$r & $p$ & acc & $n_a$ & tok & r$\to$w & w$\to$r & $p$ \\
\midrule
clean & 91.4 & 500 & 1582 & -- & -- & -- & 95.2 & 496 & 2927 & -- & -- & -- & 84.7 & 485 & 1172 & -- & -- & -- \\
typo25\_real10 & 85.5 & 497 & 1912 & 48 & 19 & 6.2e-04 & 93.9 & 495 & 2945 & 12 & 7 & 0.359 & 86.5 & 487 & 1222 & 24 & 31 & 0.418 \\
typo25\_real40 & 87.0 & 500 & 1814 & 40 & 18 & 0.006 & 93.4 & 498 & 3111 & 15 & 9 & 0.307 & 85.6 & 486 & 1246 & 28 & 33 & 0.609 \\
typo25\_real70 & 86.7 & 497 & 1843 & 43 & 19 & 0.003 & 93.9 & 494 & 2992 & 15 & 9 & 0.307 & 82.0 & 489 & 1234 & 42 & 26 & 0.069 \\
typo50\_real10 & 79.6 & 500 & 2024 & 78 & 19 & $<$1e-6 & 92.8 & 498 & 2939 & 18 & 9 & 0.124 & 81.9 & 487 & 1298 & 41 & 27 & 0.115 \\
typo50\_real40 & 79.0 & 500 & 2084 & 83 & 21 & $<$1e-6 & 91.4 & 498 & 3085 & 22 & 6 & 0.005 & 82.5 & 485 & 1357 & 42 & 28 & 0.120 \\
typo50\_real70 & 72.8 & 500 & 2094 & 108 & 15 & $<$1e-6 & 91.5 & 496 & 3087 & 25 & 9 & 0.010 & 82.4 & 484 & 1329 & 45 & 29 & 0.081 \\
typo75\_real10 & 73.8 & 488 & 2084 & 102 & 14 & $<$1e-6 & 89.1 & 494 & 2945 & 34 & 5 & 7.3e-06 & 80.7 & 488 & 1381 & 55 & 36 & 0.059 \\
typo75\_real40 & 67.3 & 495 & 2594 & 136 & 17 & $<$1e-6 & 88.1 & 496 & 3125 & 39 & 7 & 4.9e-06 & 77.3 & 485 & 1413 & 63 & 29 & 5.8e-04 \\
typo75\_real70 & 60.4 & 480 & 2306 & 164 & 14 & $<$1e-6 & 86.9 & 497 & 3138 & 46 & 7 & $<$1e-6 & 74.7 & 486 & 1453 & 75 & 26 & 1.8e-06 \\
\bottomrule
\end{tabular}}
\caption{Every configuration on every dataset. \emph{acc} is answered-only
accuracy (\%), $n_a$ the number of questions that produced a final answer,
\emph{tok} the mean generated tokens. r$\to$w and w$\to$r are paired
right-to-wrong and wrong-to-right flips against the clean condition over jointly
answered questions, with the continuity-corrected McNemar $p$-value
\DIFaddbeginFL \DIFaddBFL{(uncorrected for multiple comparisons)}\DIFaddendFL .}
\label{tab:grid}
\end{table}


\DIFaddbegin \section{\DIFaddB{Per-typo logistic fit}}
\label{app:logit}

\DIFaddend \DIFdelbegin %DIFDELCMD < \begin{table}[H]
%DIFDELCMD < %%%
\DIFdelendFL \DIFaddbeginFL \begin{table}[H]
\DIFaddendFL \centering\DIFdelbeginFL %DIFDELCMD < \scriptsize
%DIFDELCMD < %%%
\DIFdelendFL \DIFaddbeginFL \small
\DIFaddendFL \adjustbox{max width=\linewidth}{\begin{tabular}{lrrr}
\toprule
& \textbf{GSM8K} & \textbf{MATH-500} & \textbf{ARC-C} \\
\midrule
traces in the fit, $n$ & 4457 & 4466 & 4377 \\
intercept & +2.097 & +3.019 & +1.670 \\
coefficient, real-word & -0.0761 & -0.0761 & -0.0418 \\
coefficient, non-word & -0.0398 & -0.0588 & -0.0078 \\
odds multiplier, real-word & 0.927 & 0.927 & 0.959 \\
odds multiplier, non-word & 0.961 & 0.943 & 0.992 \\
odds lost per real-word typo & 7.3\% & 7.3\% & 4.1\% \\
odds lost per non-word typo & 3.9\% & 5.7\% & 0.8\% \\
\midrule
\textbf{relative loss, real / non-word} & \textbf{1.9$\times$} & \textbf{1.3$\times$} & \textbf{5.3$\times$} \\
\bottomrule
\end{tabular}}
\caption{Per-typo logistic fit behind \DIFdelbeginFL \DIFdelFL{Section~5.2}\DIFdelendFL \DIFaddbeginFL \DIFaddBFL{\S\ref{sec:realword}}\DIFaddendFL . One regression per
dataset, $\text{correct} \sim \text{num\_real} + \text{num\_nonword}$, fitted
over all answered typo traces pooled across the nine configurations. The
predictors are the counts of each kind of typo actually applied to that
question. An odds multiplier is $\exp(\text{coefficient})$: the factor by which
one additional typo of that kind multiplies the odds of a correct answer.}
\label{tab:realword-logit}
\end{table}


\DIFaddbegin \section{\DIFaddB{Paired real-word contrasts}}
\label{app:paired}

\DIFaddend \DIFdelbegin %DIFDELCMD < \begin{table}[H]
%DIFDELCMD < %%%
\DIFdelendFL \DIFaddbeginFL \begin{table}[H]
\DIFaddendFL \centering\DIFdelbeginFL %DIFDELCMD < \scriptsize
%DIFDELCMD < %%%
\DIFdelendFL \DIFaddbeginFL \small
\DIFaddendFL \adjustbox{max width=\linewidth}{\begin{tabular}{ll rrr rrr rrr}
\toprule
& & \multicolumn{3}{c}{\textbf{GSM8K}} & \multicolumn{3}{c}{\textbf{MATH-500}} & \multicolumn{3}{c}{\textbf{ARC-Challenge}} \\
\cmidrule(lr){3-5}\cmidrule(lr){6-8}\cmidrule(lr){9-11}
\textbf{$r$} & \textbf{$\rho$} & $\Delta$acc & r/w & $p$ & $\Delta$acc & r/w & $p$ & $\Delta$acc & r/w & $p$ \\
\midrule
25 & 10$\to$70 & +1.0 & 31/36 & 0.625 & +0.0 & 10/10 & 0.823 & -3.8$^{*}$ & 46/28 & 0.048 \\
25 & 10$\to$40 & +1.4 & 28/35 & 0.450 & +0.0 & 10/10 & 0.823 & -0.4 & 32/30 & 0.899 \\
25 & 40$\to$70 & -0.4 & 30/28 & 0.896 & -0.2 & 13/12 & 1.000 & -3.8$^{*}$ & 40/22 & 0.031 \\
50 & 10$\to$70 & -6.8$^{*}$ & 76/42 & 0.002 & -1.4 & 19/12 & 0.281 & +0.6 & 38/41 & 0.822 \\
50 & 10$\to$40 & -0.6 & 45/42 & 0.830 & -1.4 & 17/10 & 0.248 & +0.0 & 35/35 & 0.905 \\
50 & 40$\to$70 & -6.2$^{*}$ & 70/39 & 0.004 & -0.2 & 13/12 & 1.000 & -0.9 & 39/35 & 0.727 \\
75 & 10$\to$70 & -14.1$^{*}$ & 109/43 & $<$1e-6 & -1.8 & 28/19 & 0.243 & -6.7$^{*}$ & 60/28 & 9.5e-04 \\
75 & 10$\to$40 & -5.8$^{*}$ & 85/57 & 0.023 & -0.2 & 20/19 & 1.000 & -4.0 & 54/35 & 0.056 \\
75 & 40$\to$70 & -7.8$^{*}$ & 93/56 & 0.003 & -1.6 & 25/17 & 0.280 & -3.4 & 60/44 & 0.141 \\
\bottomrule
\end{tabular}}
\caption{Controlled paired contrast behind \DIFdelbeginFL \DIFdelFL{Section~5.2}\DIFdelendFL \DIFaddbeginFL \DIFaddBFL{\S\ref{sec:realword}}\DIFaddendFL , at every typo
rate and every pair of real-word ratios. Within a fixed $r$ the same words of
the same questions are corrupted in the same positions, so only the \emph{kind}
of corruption changes. $\Delta$acc is the accuracy of the higher-$\rho$ variant
minus the lower, over questions answered in both; \emph{r/w} is right-to-wrong
against wrong-to-right flips; $p$ is a continuity-corrected McNemar test and
$^{*}$ marks $p<0.05$ \DIFaddbeginFL \DIFaddBFL{before multiple-comparison correction}\DIFaddendFL .}
\label{tab:realword-paired}
\end{table}


\DIFaddbegin \section{\DIFaddB{Reasoning length and self-doubt markers}}
\label{app:support}

\DIFaddend \DIFdelbegin %DIFDELCMD < \begin{table}[H]
%DIFDELCMD < %%%
\DIFdelendFL \DIFaddbeginFL \begin{table}[H]
\DIFaddendFL \centering\scriptsize\DIFaddbeginFL \setlength{\tabcolsep}{2.6pt}
\DIFaddendFL \adjustbox{max width=\linewidth}{\begin{tabular}{l rrrrrr rrrrrr rrrrrr}
\toprule
& \multicolumn{6}{c}{\textbf{GSM8K}} & \multicolumn{6}{c}{\textbf{MATH-500}} & \multicolumn{6}{c}{\textbf{ARC-Challenge}} \\
\cmidrule(lr){2-7}\cmidrule(lr){8-13}\cmidrule(lr){14-19}
\textbf{config} & tok & tok$_{c}$ & tok$_{w}$ & uncert. & 2nd-g. & u.w/c & tok & tok$_{c}$ & tok$_{w}$ & uncert. & 2nd-g. & u.w/c & tok & tok$_{c}$ & tok$_{w}$ & uncert. & 2nd-g. & u.w/c \\
\midrule
clean & 1582 & 1406 & 3456 & 7.4 & 13.0 & 1.72 & 2927 & 2803 & 5352 & 6.8 & 12.6 & 1.70 & 1172 & 1025 & 1989 & 14.0 & 11.3 & 1.66 \\
typo25\_real10 & 1912 & 1511 & 4281 & 12.3 & 14.7 & 1.83 & 2945 & 2826 & 4783 & 7.0 & 11.8 & 1.63 & 1222 & 1120 & 1870 & 14.5 & 11.8 & 1.57 \\
typo25\_real40 & 1814 & 1567 & 3470 & 12.6 & 14.0 & 1.90 & 3111 & 2900 & 6075 & 7.5 & 13.7 & 1.52 & 1246 & 1114 & 2026 & 13.4 & 11.0 & 1.71 \\
typo25\_real70 & 1843 & 1537 & 3843 & 13.2 & 14.5 & 2.12 & 2992 & 2859 & 5047 & 8.2 & 13.2 & 1.68 & 1234 & 1066 & 1998 & 14.6 & 10.8 & 1.94 \\
typo50\_real10 & 2024 & 1554 & 3856 & 15.4 & 14.6 & 1.87 & 2939 & 2763 & 5191 & 7.9 & 12.7 & 1.97 & 1298 & 1184 & 1814 & 14.7 & 11.0 & 1.50 \\
typo50\_real40 & 2084 & 1631 & 3787 & 15.9 & 14.7 & 1.63 & 3085 & 2839 & 5694 & 8.8 & 12.7 & 1.96 & 1357 & 1206 & 2069 & 16.3 & 11.7 & 1.92 \\
typo50\_real70 & 2094 & 1501 & 3682 & 17.1 & 14.8 & 1.65 & 3087 & 2854 & 5600 & 8.5 & 13.1 & 1.90 & 1329 & 1180 & 2032 & 16.1 & 11.1 & 1.67 \\
typo75\_real10 & 2084 & 1588 & 3479 & 18.7 & 14.4 & 1.62 & 2945 & 2787 & 4236 & 8.1 & 12.7 & 2.06 & 1381 & 1232 & 2006 & 17.1 & 11.5 & 1.48 \\
typo75\_real40 & 2594 & 1786 & 4253 & 21.3 & 14.7 & 1.60 & 3125 & 2881 & 4932 & 9.0 & 12.7 & 2.06 & 1413 & 1263 & 1925 & 18.4 & 11.3 & 1.38 \\
typo75\_real70 & 2306 & 1652 & 3306 & 21.9 & 14.1 & 1.55 & 3138 & 2912 & 4640 & 9.3 & 12.5 & 1.75 & 1453 & 1263 & 2015 & 19.2 & 11.6 & 1.47 \\
\bottomrule
\end{tabular}}
\caption{Supporting numbers for \DIFdelbeginFL \DIFdelFL{Sections~5.1 }\DIFdelendFL \DIFaddbeginFL \DIFaddBFL{\S\ref{sec:acc} }\DIFaddendFL and \DIFdelbeginFL \DIFdelFL{~5.4}\DIFdelendFL \DIFaddbeginFL \DIFaddBFL{\S\ref{sec:doubt}}\DIFaddendFL .
\emph{tok} is mean generated tokens over all answered traces; \emph{tok}$_{c}$
and \emph{tok}$_{w}$ are the same mean \DIFdelbeginFL \DIFdelFL{taken }\DIFdelendFL over traces that ended in a correct and
in a wrong answer. \emph{uncert.} and \emph{2nd-g.}\ are uncertainty and
second-guessing markers per 1{,}000 reasoning words, and \emph{u.w/c} is
uncertainty density in wrong-answer traces divided by correct-answer traces.
Traces \DIFaddbeginFL \DIFaddBFL{generally }\DIFaddendFL lengthen with the typo rate \DIFdelbeginFL \DIFdelFL{on every dataset}\DIFdelendFL \DIFaddbeginFL \DIFaddGFL{(MATH-500 at $\rho{=}10$ is flat)}\DIFaddendFL ,
a wrong answer is preceded by a longer trace than a correct one in every
configuration, and uncertainty climbs while second-guessing stays \DIFaddbeginFL \DIFaddBFL{roughly
}\DIFaddendFL flat.}
\label{tab:support}
\end{table}
\DIFaddbegin 


\section{\DIFaddB{Word-level repair behaviour}}
\label{app:repair}
\DIFaddend 

\DIFdelbegin %DIFDELCMD < \begin{table}[H]
%DIFDELCMD < %%%
\DIFdelendFL \DIFaddbeginFL \begin{table}[H]
\DIFaddendFL \centering\scriptsize\DIFaddbeginFL \setlength{\tabcolsep}{2.6pt}
\DIFaddendFL \adjustbox{max width=\linewidth}{\begin{tabular}{ll rrrrrr rrrrrr rrrrrr}
\toprule
& & \multicolumn{6}{c}{\textbf{GSM8K}} & \multicolumn{6}{c}{\textbf{MATH-500}} & \multicolumn{6}{c}{\textbf{ARC-Challenge}} \\
\cmidrule(lr){3-8}\cmidrule(lr){9-14}\cmidrule(lr){15-20}
\textbf{$r$} & \textbf{$\rho$} & silent & flag & mis. & unused & mis$_{c}$ & mis$_{w}$ & silent & flag & mis. & unused & mis$_{c}$ & mis$_{w}$ & silent & flag & mis. & unused & mis$_{c}$ & mis$_{w}$ \\
\midrule
25 & 10 & 73.3 & 11.9 & 2.9 & 12.0 & 2.5 & 5.0 & 78.0 & 2.3 & 1.2 & 18.4 & 1.3 & 0.4 & 80.0 & 5.8 & 2.1 & 12.1 & 1.8 & 3.9 \\
25 & 40 & 70.2 & 15.1 & 4.3 & 10.4 & 3.9 & 6.8 & 73.1 & 5.8 & 1.6 & 19.5 & 1.3 & 3.9 & 74.8 & 9.1 & 3.1 & 13.1 & 2.7 & 5.2 \\
25 & 70 & 64.6 & 19.3 & 4.8 & 11.3 & 4.3 & 7.7 & 71.4 & 10.2 & 1.7 & 16.8 & 1.4 & 4.1 & 69.6 & 12.2 & 5.3 & 12.9 & 5.0 & 6.4 \\
50 & 10 & 65.8 & 18.7 & 4.7 & 10.9 & 3.9 & 7.7 & 75.0 & 4.3 & 1.4 & 19.4 & 1.3 & 2.5 & 76.9 & 8.1 & 2.5 & 12.5 & 2.0 & 4.9 \\
50 & 40 & 60.6 & 23.4 & 5.2 & 10.8 & 5.1 & 5.8 & 70.2 & 8.3 & 2.3 & 19.2 & 2.2 & 3.2 & 71.8 & 11.7 & 3.7 & 12.7 & 3.4 & 5.1 \\
50 & 70 & 55.9 & 25.3 & 6.6 & 12.2 & 6.1 & 7.8 & 67.2 & 11.3 & 2.9 & 18.6 & 2.4 & 6.0 & 66.7 & 14.2 & 5.4 & 13.7 & 5.1 & 6.2 \\
75 & 10 & 55.8 & 26.1 & 6.9 & 11.3 & 6.1 & 8.7 & 70.4 & 8.2 & 2.1 & 19.3 & 1.7 & 3.7 & 71.0 & 13.4 & 3.4 & 12.2 & 2.9 & 5.2 \\
75 & 40 & 47.4 & 32.9 & 8.3 & 11.4 & 6.8 & 11.1 & 68.4 & 9.6 & 2.4 & 19.5 & 1.9 & 4.8 & 65.8 & 16.3 & 5.0 & 12.8 & 4.7 & 6.2 \\
75 & 70 & 46.2 & 30.9 & 9.6 & 13.3 & 8.0 & 11.9 & 62.2 & 15.2 & 4.0 & 18.5 & 3.5 & 6.0 & 59.4 & 20.3 & 6.9 & 13.3 & 6.1 & 9.3 \\
\bottomrule
\end{tabular}}
\caption{Every corrupted word in every configuration, classified by how the
reasoning trace handles it, as a percentage of the corrupted words in that
configuration. \emph{silent} / \emph{flag} / \emph{mis.} / \emph{unused} are the
four categories of \DIFdelbeginFL \DIFdelFL{Section~5.3}\DIFdelendFL \DIFaddbeginFL \DIFaddBFL{\S\ref{sec:repair}}\DIFaddendFL . \emph{mis}$_{c}$ and \emph{mis}$_{w}$ are
the misread rate among questions the model answered correctly and among those it
answered wrongly, which is the comparison behind the claim that misreading is
associated with failure. Figure~\ref{fig:repair} shows the same data pooled
along each axis.}
\label{tab:repair}
\end{table}


\end{document}
